\documentclass[a4paper,11pt]{article}

% ============================================================
%  QUANT INTERVIEW PREP TEMPLATE
%  Adapted from a class-notes template.
%
%  Three sections:
%    1) Math / Probability Theory   -> theorem/definition/proof + key-result boxes
%    2) Brainteasers                -> numbered problems + toggleable solutions
%    3) Code                        -> paradigm notes + syntax-highlighted practice
%
%  SOLUTION TOGGLE (self-testing): flip the switch below.
%    \showsolutionstrue   -> solutions printed (review / check mode)
%    \showsolutionsfalse  -> solutions hidden  (quiz yourself, then recompile)
% ============================================================

\usepackage[english]{babel}
\usepackage[utf8]{inputenc}       % utf8x is deprecated; utf8 is the modern default
\usepackage[T1]{fontenc}
\usepackage{lmodern}              % scalable fonts for all shapes (incl. italic monospace)
\usepackage[a4paper,margin=1in]{geometry}

\usepackage{amsmath}
\usepackage{amssymb}
\usepackage{amsthm}
\usepackage{mathtools}            % superset of amsmath (\coloneqq, \DeclarePairedDelimiter, etc.)
\usepackage{graphicx}
\usepackage{xcolor}
\usepackage{listings}
\usepackage[most]{tcolorbox}
\usepackage{comment}             % used to drop solution blocks when hidden
\usepackage{parskip}             % paragraph spacing instead of indent, reads better for notes
\usepackage{enumitem}
\usepackage[colorlinks=true,linkcolor=blue!55!black,urlcolor=blue!55!black,citecolor=blue!55!black]{hyperref}
\setlength{\marginparwidth}{2cm}            % avoids a todonotes load-time warning
\usepackage[colorinlistoftodos]{todonotes}  % \todo{revisit this} to flag weak spots
\setuptodonotes{inline}                     % full-width inline banners, never clip off the page

% ---------- theorem-like environments ----------
\theoremstyle{plain}
\newtheorem{theorem}{Theorem}[section]
\newtheorem{lemma}[theorem]{Lemma}
\newtheorem{corollary}[theorem]{Corollary}
\newtheorem{proposition}[theorem]{Proposition}
\theoremstyle{definition}
\newtheorem{definition}[theorem]{Definition}
\newtheorem{example}[theorem]{Example}
\theoremstyle{remark}
\newtheorem*{remark}{Remark}

% ---------- probability / math shorthand ----------
\newcommand{\R}{\mathbb{R}}
\newcommand{\N}{\mathbb{N}}
\newcommand{\Z}{\mathbb{Z}}
\newcommand{\E}{\mathbb{E}}
\newcommand{\Prob}{\mathbb{P}}
\newcommand{\Var}{\operatorname{Var}}
\newcommand{\Cov}{\operatorname{Cov}}
\newcommand{\ind}{\mathbf{1}}                   % indicator function (no extra font package)
\newcommand{\given}{\,\vert\,}
\newcommand{\indep}{\perp\!\!\!\perp}
\DeclareMathOperator*{\argmax}{arg\,max}
\DeclareMathOperator*{\argmin}{arg\,min}

% ---------- key-result highlight box (math section) ----------
\newtcolorbox{keybox}[1][Key result]{%
  breakable, colback=blue!4, colframe=blue!55!black, coltitle=white,
  fonttitle=\bfseries, title=#1, boxrule=0.6pt, arc=2pt,
  left=6pt, right=6pt, top=4pt, bottom=4pt}

% ---------- numbered problem environment (brainteasers) ----------
% Optional argument holds metadata, e.g. \begin{problem}[Brainstellar, Medium]
\makeatletter
\newcounter{problem}[section]
\renewcommand{\theproblem}{\thesection.\arabic{problem}}
\newenvironment{problem}[1][]{%
  \par\medskip\refstepcounter{problem}%
  \noindent\textbf{Problem~\theproblem}%
  \def\prob@meta{#1}%
  \ifx\prob@meta\@empty\else\ \textnormal{\itshape[#1]}\fi%
  \par\nobreak\smallskip
}{\par\medskip}
\makeatother

% ---------- toggleable solution environment ----------
\newif\ifshowsolutions
\showsolutionstrue        % <-- flip to \showsolutionsfalse to hide all solutions
\ifshowsolutions
  \newenvironment{solution}%
    {\par\smallskip\noindent\textbf{Solution.}\enspace}%
    {\par\medskip}
\else
  \excludecomment{solution}
\fi

% ---------- code listing styles ----------
\definecolor{lstbg}{RGB}{248,248,248}
\definecolor{lstkw}{RGB}{0,0,170}
\definecolor{lstcm}{RGB}{95,120,95}
\definecolor{lstst}{RGB}{160,60,60}
\definecolor{cream}{RGB}{253,251,245}
\lstdefinestyle{common}{%
  backgroundcolor=\color{lstbg}, basicstyle=\ttfamily\small,
  keywordstyle=\color{lstkw}\bfseries, commentstyle=\color{lstcm}\itshape,
  stringstyle=\color{lstst}, numbers=left, numberstyle=\tiny\color{gray},
  numbersep=8pt, showstringspaces=false, breaklines=true, frame=single,
  framerule=0.3pt, rulecolor=\color{gray!40}, tabsize=4, captionpos=b,
  xleftmargin=16pt, aboveskip=1em, belowskip=1em}
\lstdefinestyle{python}{style=common, language=Python}
\lstdefinestyle{sql}{style=common, language=SQL}
\lstset{style=common}

% \complexity{time}{space}  -> "Complexity: time O(...), space O(...)."
\newcommand{\complexity}[2]{\noindent\textbf{Complexity:}\enspace time $#1$, space $#2$.\par}

% ============================================================

\title{HFT Ops Interview Preparation \\[0.4em]
    \large \textit{Probability \, / \, Brainteasers \, / \, Code}}
\author{Gabriel Kinshuk}
\date{\today}

\begin{document}
\pagecolor{cream}
\maketitle
\tableofcontents
\newpage

% ============================================================
\section{Math / Probability Theory}
% ============================================================

\subsection{Combinatorics}

A counting problem is pinned down by two questions. Does the order of the chosen elements matter, and may an element be chosen more than once. The formulas below are organized around those two axes, plus a correction for identical elements and two escape hatches (the complement and inclusion-exclusion) for when a direct count fans out into many cases.

\subsubsection{Permutations}

\begin{definition}[Permutation]
A permutation is an ordered selection. Choosing $k$ items from $n$ distinct options where order matters gives
\[
P(n, k) = \frac{n!}{(n-k)!}.
\]
\end{definition}

The formula counts position by position. The first slot has $n$ candidates, the second has $n-1$ once one item is placed, and so on for $k$ slots, giving the falling product $n(n-1)\cdots(n-k+1)$. Writing that product as $\frac{n!}{(n-k)!}$ just cancels the unused tail $(n-k)!$.

\begin{definition}[Permutations with identical elements]
When the $n$ items being arranged include groups of identical copies with frequencies $n_1, n_2, \dots, n_k$, the number of distinguishable arrangements is
\[
\frac{n!}{n_1!\, n_2! \cdots n_k!}.
\]
\end{definition}

Start from the $n!$ arrangements that would exist if every item were distinct. Each identical group of size $n_i$ has $n_i!$ internal orderings that all produce the same visible arrangement, so dividing by each $n_i!$ removes exactly those duplicates.

\subsubsection{Combinations}

\begin{definition}[Combination]
A combination is an unordered selection. Choosing $k$ items from $n$ distinct options where order does not matter gives
\[
C(n,k) = \binom{n}{k} = \frac{n!}{k!\,(n-k)!}.
\]
\end{definition}

Combinations come from permutations by discarding the ordering. Each unordered set of $k$ items can be arranged in $k!$ ways, and $P(n,k)$ counts all of those as distinct, so
\[
C(n,k) = \frac{P(n,k)}{k!}.
\]

\subsubsection{Distributing Indistinguishable Objects (Stars and Bars)}

\begin{theorem}[Stars and bars]
The number of ways to place $n$ identical objects into $k$ distinct bins, where a bin may hold any number of objects including zero, is
\[
\binom{n + k - 1}{\,k - 1\,} = C(n + k - 1,\, k - 1).
\]
\end{theorem}

\begin{proof}
Lay the $n$ objects in a row as stars and insert $k-1$ dividers (bars) to cut the row into $k$ groups. Reading left to right, objects before the first bar fall in bin 1, objects between the first and second bar fall in bin 2, and so on through bin $k$. Every placement of the bars produces one distribution, and every distribution corresponds to exactly one placement. The row has $n + (k-1)$ positions, and a distribution is fixed by choosing which $k-1$ of them hold bars, which is $\binom{n+k-1}{k-1}$.
\end{proof}

\subsubsection{Inclusion-Exclusion}

\begin{theorem}[Inclusion-Exclusion]
The size of a union is the alternating sum of the sizes of all intersections. For two and three sets,
\[
|A \cup B| = |A| + |B| - |A \cap B|,
\]
\[
|A \cup B \cup C| = |A| + |B| + |C| - |A\cap B| - |A\cap C| - |B\cap C| + |A\cap B\cap C|.
\]
\end{theorem}

The mechanism is bookkeeping on how many times each element gets counted. Summing the individual set sizes counts an element that lies in two sets twice, so subtracting the pairwise intersections corrects it back to one. An element in all three sets is now counted $3 - 3 = 0$ times, so adding the triple intersection restores it to one. The pattern continues to any number of sets: terms over an odd number of sets are added and terms over an even number are subtracted, which leaves every element counted exactly once no matter how many sets contain it.

\begin{keybox}[Which counting tool]
\begin{center}
\begin{tabular}{lll}
\toprule
\textbf{Order matters?} & \textbf{Repeats allowed?} & \textbf{Count} \\
\midrule
yes & no  & $P(n,k) = \dfrac{n!}{(n-k)!}$ \\[0.6em]
no  & no  & $C(n,k) = \dbinom{n}{k}$ \\[0.6em]
yes & yes & $n^k$ \\[0.3em]
no  & yes & $\dbinom{n+k-1}{k-1}$ \; (stars and bars) \\
\bottomrule
\end{tabular}
\end{center}
Two escape hatches handle counts that fan out into many subcases. Use the \textbf{complement} ($1 - P$ of the unwanted case) on ``at least one'' phrasing, and use \textbf{inclusion-exclusion} on ``divisible by $a$ or $b$'' style overlaps.
\end{keybox}

\newpage

\subsubsection{Select Worked Problems}

\begin{problem}
How many ways can you arrange the letters in the word BANANA?
\end{problem}

\begin{solution}
There are $n = 6$ elements. Repeated items have frequency 3 and 2 (Ns and As). Then the number of arrangements are $\frac{6!}{3! \times 2!} = \frac{720}{12} = 60$.
\end{solution}

\medskip

\begin{problem}
You roll 2 dice. How many distinct outcomes are there? What if you can't tell the dice apart?
\end{problem}

\begin{solution}
With two dice, there are $6 \times 6 = 36$ outcomes if we can tell the dice apart. If we \textit{can't} tell the dice apart, then we can solve either logically or formulaically:

Logically, we can see that from the original 36 combinations, there are 6 outcomes where the dice are the same, leaving 30 outcomes using mixed pairs of dice. Of these, half are duplicate when we can't tell the dice apart, leaving 15 outcomes. We can then add the 6 double outcomes back in for a total of 21 outcomes.

Formulaically, we can solve using a similar approach, taking the sum of mixed pair outcomes and double outcomes as $C(6, 2) + C(6, 1) = 15 + 6 = 21$.

Lastly, we can consider this problem using the stars and bars approach, where we notice that there are $n=2$ identical dice and $6$ distinct bins with $k=5$ dividers between them. Then we can count the outcomes by $C(2 + 6 - 1, 6 -1) = C(7, 5) = 21$
\end{solution}

\medskip

\begin{problem}
How many ways can you distribute 8 identical balls into 3 distinct boxes?
\end{problem}

\begin{solution}
This is a direct application of stars and bars with $n=8$ and $k=3$. The number of outcomes is $C(8+3-1, 3-1) = C(10,2) = \frac{10 \times 9 \times 8 \times \cdots \times 3}{2} = 45$.
\end{solution}

\medskip

\begin{problem}
A deck has 52 cards. How many 5-card hands contain at least one ace?
\end{problem}

\begin{solution}
Since this solution asks for the number of hands containing at least one ace, it will be simpler to compute the complement. This is the number of hands containing no aces, we we then subtract from the total number of 5 card hands: $C(52, 5) - C(48, 5) = 886,656$.
\end{solution}

\medskip

\begin{problem}
You have 5 books on a shelf. In how many ways can you choose 3 books and arrange them left-to-right on a new shelf?
\end{problem}

\begin{solution}
This is simply $C(5,3) \times P(3, 3) = \frac{5 \times 4 \times 3}{3!} \times \frac{3!}{0!} = \frac{60 \times 3!}{3!} = 60$.
\end{solution}

\medskip

\begin{problem}
A pizza place offers 8 toppings. How many different pizzas can you make if you can choose any number of toppings (including none)?
\end{problem}

\begin{solution}
We can solve this by taking the summation of $C(8, i)$ from $i=0$ to $i=8$:
\[
Total\ outcomes = \sum\limits_{i=0}^{8} C(8,i) = 256
\]

There is a more succinct solution we can find if we approach the problem seeing each topping as the binary choice of being either on or off of the pizza, then the total outcomes are $2^8 = 256$.
\end{solution}

\medskip

\begin{problem}
How many ways can 12 people be split into 3 groups of 4?
\end{problem}

\begin{solution}
We choose each group independently with the remaining members of the pool, and then divide by the reorderings since the groups are indistinguishable:
\[
\frac{C(12, 4) \times C(8, 4) \times C(4, 4)}{3!} = \frac{495 \times 70 \times 1}{6} = 34,650.
\]
\end{solution}

\bigskip

\subsection{Probability}

\begin{definition}[Sample Space]
The sample space $\Omega$ is the set of all possible outcomes of an experiment.
\end{definition}

Most problems begin by fixing the sample space and then identifying the subset of outcomes that count as favourable. Getting the sample space right, which means choosing the correct unit of observation and the correct notion of ``equally likely'', settles most of the difficulty before any arithmetic starts.

\begin{definition}[Event]
An event is a subset of the sample space, the collection of outcomes of interest. Its probability is the total weight of the outcomes it contains.
\end{definition}

\subsubsection{Complement}

\begin{keybox}[Complement]
For any event $A$,
\[
P(A) = 1 - P(A^c),
\]
where $A^c$ is the event that $A$ does not occur.
\end{keybox}

The complement is the first move whenever the direct event splits into many cases while its negation is a single clean case. ``At least one'' is the signature cue: the direct count runs over one, two, three or more successes, while the negation is the single event of zero successes. The birthday and ``at least one head'' problems below both collapse this way.

\subsubsection{Independence}

\begin{definition}[Independence]
Two events are independent when the occurrence of one leaves the probability of the other unchanged. Formally,
\[
P(A \cap B) = P(A)\,P(B).
\]
\end{definition}

The multiplication rule is the practical test. Drawing \textit{with} replacement keeps every draw independent because the pool is restored each time. Drawing without replacement makes the draws dependent, since removing an item changes both the size and the composition of the pool for the next draw. Independence is a statement about probabilities and is a separate idea from mutual exclusivity, which is a statement about outcomes.

\subsubsection{Mutual Exclusivity}

\begin{definition}[Mutual exclusivity]
Two events are mutually exclusive when they cannot both occur, so $A \cap B = \emptyset$. For such events the union probability adds:
\[
P(A \cup B) = P(A) + P(B).
\]
\end{definition}

This is the special case of inclusion-exclusion where the intersection is empty. The general identity $P(A \cup B) = P(A) + P(B) - P(A \cap B)$ carries the overlap term $P(A \cap B)$, and mutual exclusivity sets that term to zero. A common error is to add probabilities for events that can co-occur, which double counts the overlap.


\subsubsection{Select Worked Problems}

\begin{problem}
How many people do you need in a room for a >50\% chance that two share a birthday?
\end{problem}

\begin{solution}
Computing this directly is impractical, since there are many ways for multiple people to share a birthday. Instead, using the complement, we simply subtract the probability that \textit{no} people share a birthday from 1:
\[
P(\text{all different}) = \prod\limits_{i=1}^{n} \frac{365 - i + 1}{365}
\]
We multiply additional factors until the quantity drops \textit{below} 0.5 (because then $1 - P$ is greater than 0.5). This turns out to be $n=23$, where $P(\text{all different}) \approx 0.493$ and thus $P(\text{birthday is shared}) \approx 0.507$.
\end{solution}

\begin{problem}
You flip a fair coin 3 times. What's the probability of getting exactly 2 heads?
\end{problem}

\begin{solution}
Flipping a coin three times has eight outcomes. Of these, three involve exactly two heads, thus the probability is $\frac{3}{8} = 0.375$.

If we don't want to enumerate each of the eight options and count directly, we could use $C(3, 2) = 3$ to find the number of outcomes involving two heads.
\end{solution}


\begin{problem}
Roll two dice. What's the probability the sum is 7?
\end{problem}

\begin{solution}
There are 36 outcomes of rolling two dice. Of these, the pairs that equal 7 are $(1,6), (2,5), (3,4), (4,3), (5,2), (6,1)$. Then the probability is $\frac{6}{36} \approx 0.167$.
\end{solution}

\begin{problem}
You draw 2 cards from a standard 52-card deck (without replacement). What's the probability both are hearts?
\end{problem}

\begin{solution}
Without replacement, the draws are dependent. Then, the probability is:
\[
\frac{13}{52} \times \frac{12}{51} \approx 0.0588
\]
Note that \textit{both} the numerator and the denominator shrank after we removed the first card.

We could also solve this by counting, if we divide the number of ways to pick two hearts by the total number of ways to pick two cards:
\[
\frac{C(13,2)}{C(52,2)} = \frac{78}{1326} \approx 0.0588
\]
\end{solution}

\begin{problem}
You flip a fair coin 5 times. What's the probability of getting at least one head?
\end{problem}

\begin{solution}
We can use the complement, the probability of getting \textit{no} heads, subtracted from 1:
\[
1 - \frac{1}{2^5} = 0.96875
\]
\end{solution}

\begin{problem}
You roll a fair die twice. What's the probability the second roll is strictly greater than the first?
\end{problem}

\begin{solution}
THere are 36 total ways to roll two dice. If the first roll is a 1, then there are 5 ways to roll higher on the second die. Similarly, there are 4 ways to roll higher if we rolled a 2, and so on.

Thus the number of favourable outcomes $1 + 2 + \cdots + 5$ or the triangle number of 5, $T_5 = 15$. The final answer is:
\[
\frac{15}{36} \approx 0.417
\]

Alternatively, we could solve by symmetry, calling roll 1 and 2 as $A$ and $B$, respectively. If we notice that $P(A > B) = P(B < A)$, and $P(A = B) = \frac{6}{36}$, then we evenly divide the remaining outcomes by 2 to get $\frac{15}{36}$.
\end{solution}

\begin{problem}
A box has 5 white and 3 black balls. You draw 3 without replacement. What's the probability of drawing exactly 2 white and 1 black?
\end{problem}

\begin{solution}
The total ways to draw 3 balls are $C(8, 3) = 56$. We then find the ways to choose 2 white balls and 1 black ball individually, which is $C(5,2)$ and $C(3, 1)$, then the answer is:
\[
\frac{C(5,2)\times C(3,1)}{C(8,3)} = \frac{10 \times 3}{56} \approx 0.536
\]
Since we remove the balls without replacement, this problem follows the hypergeometric distribution.
\end{solution}

\begin{problem}
You randomly pick a number from 1 to 100. What's the probability it's divisible by 3 or 5?
\end{problem}

\begin{solution}
There are 33 numbers divisible by 3, and 20 numbers divisible by 5. The numbers divisible by both are the multiples of 15 of which there are 6 (15, 30, 45, 60, 75, 90). Thus:
\[
P_{Div}(A \cup B) = \frac{33}{100} + \frac{20}{100} - \frac{6}{100} = 0.47
\]
\end{solution}

\begin{problem}
A family has 3 children. Assuming each child is equally likely to be a boy or girl, what's the probability that all 3 are the same gender?
\end{problem}

\begin{solution}
There are eight outcomes, of those there are two in which all children are the same gender, so the answer is $\frac{2}{8} = 0.25$.
\end{solution}


\begin{problem}
You throw 3 dice. What's the probability ALL three show different numbers?
\end{problem}

\begin{solution}
There are $6^3 = 216$ total ways to throw three dice. After we throw the first dice, there are 5 options remaining, and 4 after the second dice. Then the answer is:
\[
\frac{6 \times 5 \times 4}{6^3} = \frac{120}{216} \approx 0.556
\]
\end{solution}

\begin{problem}
You toss a fair coin until you get heads. What's the probability you need exactly 3 tosses?
\end{problem}

\begin{solution}
We stop at the first heads, so needing exactly 3 tosses means the first two tosses are tails and the third is heads: T, T, H. Since tosses are independent and each outcome has probability $\frac{1}{2}$,
\[
P(\text{exactly 3}) = \tfrac{1}{2} \times \tfrac{1}{2} \times \tfrac{1}{2} = \tfrac{1}{8}.
\]
We do not enumerate anything else, because exactly one sequence produces this event.
\end{solution}


\begin{remark}[Two ways probability gets counted]
The reason we can pin down this single case without listing every outcome is that this problem lives in a different kind of sample space than counting problems like the birthday question.

When every outcome is equally likely, probability is combinatorial. You divide favorable outcomes by total outcomes, so you have to characterize the entire sample space. Each outcome carries no probability on its own and earns weight only by being counted against the total.

When outcomes are not equally likely, each outcome carries its own probability directly, usually a product of independent step probabilities. Here the sequence $\text{TTH}$ has probability $\tfrac{1}{2}\cdot\tfrac{1}{2}\cdot\tfrac{1}{2} = \tfrac{1}{8}$ by itself, with no denominator to divide by. That is why one case can be solved in isolation.

These are the same axioms applied to different spaces. Counting favorable over total is the special case that appears when all per-outcome probabilities are equal, so the sum of probabilities collapses into a plain count. The coin problem is the general version, where the per-outcome probabilities differ and you take the one you want.
\end{remark}

\begin{problem}
In a room of 5 people, what's the probability that at least two share a birthday? (Assume 365 days, equally likely.)
\end{problem}

\begin{solution}
Again, we use the complement to find 1 minus the probability that no one shares a birthday:
\[
1 - \left ( \frac{365}{365} \times \frac{364}{365} \times \frac{363}{365} \times \frac{362}{365} \times \frac{361}{365}  \right ) \approx 1 - 0.9729 \approx 0.0271
\]
\end{solution}

\bigskip

\subsection{Expected Value}

\begin{definition}[Expected Value]
The expected value of a random variable is the probability-weighted average of its outcomes:
\[
\E[X] = \sum_i x_i \, P(X = x_i).
\]
It is the long-run mean over many repetitions. The expected value need not be an attainable outcome (a fair die has $\E[X] = 3.5$) and need not be the most likely one.
\end{definition}

\begin{keybox}[Linearity of expectation]
For any random variables $X$ and $Y$ and any scalars $a, b$,
\[
\E[aX + bY] = a\,\E[X] + b\,\E[Y].
\]
This holds with no independence assumption. Paired with indicator variables it is the highest-yield tool for expectation brainteasers, since it turns a hard expected count into a sum of easy probabilities.
\end{keybox}

\begin{definition}[Indicator Random Variable]
An indicator $\ind\{A\}$ equals $1$ when event $A$ occurs and $0$ otherwise. Its expectation is the event's probability,
\[
\E[\ind\{A\}] = 1 \cdot P(A) + 0 \cdot P(A^c) = P(A).
\]
\end{definition}

The variable itself is trivial. The leverage is the decomposition it enables. To find the expected \textit{count} of how many events in a family occur, write that count as a sum of indicators, one per event, then apply linearity:
\[
\E\!\left[\sum_i \ind\{A_i\}\right] = \sum_i \E[\ind\{A_i\}] = \sum_i P(A_i).
\]
Because linearity ignores dependence, this works even when the events are entangled, which is what collapses an otherwise messy count into a sum of easy probabilities. The shuffle and fixed-point problems below are the standard instances.

\begin{theorem}[Law of Total Expectation]
If the cases $B_1, \dots, B_n$ partition the sample space, meaning they are mutually exclusive and cover every outcome, then
\[
\E[X] = \sum_{i=1}^{n} \E[X \mid B_i] \cdot P(B_i).
\]
\end{theorem}

This regroups the definition of expected value by cases. Rather than averaging over individual outcomes, average the conditional mean inside each case and weight it by the case probability. Expanding each $\E[X \mid B_i]$ recovers the flat sum, so the identity adds no new information, but it changes what is computable. It is the tool of choice in three situations:
\begin{itemize}[itemsep=-1ex, label=--]
    \item you are handed conditional averages and case probabilities rather than the full outcome distribution
    \item $X$ is continuous, so the flat sum $\sum_x x \, P(X=x)$ does not apply while conditioning on discrete cases still does
    \item $\E[X \mid B_i]$ refers back to $\E[X]$ itself, which turns the conditioning into a solvable equation
\end{itemize}
The cases must be disjoint and exhaustive. Overlapping or missing cases break the weighting.

\begin{example}[Conditioning on cases]
A commute time depends on the weather. Sunny days are $60\%$ of days with a $20$ minute average, and rainy days are $40\%$ with a $35$ minute average. The two cases partition every day, so
\[
\E[\text{commute}] = 0.6(20) + 0.4(35) = 12 + 14 = 26 \text{ minutes.}
\]
\end{example}

\begin{example}[Recursive expectation]
Take the expected number of rolls of a die until the first 6, and call it $E$. One roll always happens. With probability $\frac{1}{6}$ it is a 6 and the process stops, and with probability $\frac{5}{6}$ it is not, which leaves the same problem ahead with the same expectation $E$:
\[
E = 1 + \tfrac{5}{6}E \;\longrightarrow\; \tfrac{1}{6}E = 1 \;\longrightarrow\; E = 6.
\]
This is the mean of the geometric distribution, the expected number of trials until the first success, which is $\frac{1}{p}$ with $p = \frac{1}{6}$.
\end{example}



\subsubsection{Select Worked Problems}

\begin{problem}
A game costs \$1. You flip a coin: heads you win \$3, tails you win nothing. What's $\mathbb{E}[\text{Profit}]$?
\end{problem}

\begin{solution}
$EV = 0.5 \cdot 2 + 0.5 \cdot -1 = +\ 0.5$
\end{solution}

\begin{problem}
Roll a die. If it shows 6, you roll again and add. Otherwise, your score is the number shown. What's the expected score?
\end{problem}

\begin{solution}
\[
EV = 3.5 + \frac{1}{6} \cdot EV \longrightarrow \frac{5}{6} \cdot EV = 3.5 \longrightarrow EV = 4.2
\]

\end{solution}

\begin{problem}
Shuffle a standard deck. What's the expected number of cards that are in their "correct" position (card k is in position k)?
\end{problem}

\begin{solution}
This is linearity of expectation applied to indicators. Let $I_k = 1$ if card $k$ lands in position $k$, and $0$ otherwise. In a random shuffle each card is equally likely to occupy any of the 52 positions, so $\mathbb{E}[I_k] = P(I_k = 1) = \frac{1}{52}$. The number of correctly placed cards is $M = \sum_{k=1}^{52} I_k$, and by linearity
\[
\mathbb{E}[M] = \sum_{k=1}^{52} \mathbb{E}[I_k] = 52 \cdot \frac{1}{52} = 1.
\]
The $I_k$ are dependent, but linearity does not require independence, so that never matters.
\end{solution}

\newpage

\begin{problem}
You flip a coin until you get heads. How many flips do you expect?
\end{problem}

\begin{solution}
\[
EV = 1 + 0.5 \cdot EV \longrightarrow 0.5 \cdot EV = 1 \longrightarrow EV = 2
\]
\end{solution}

\begin{problem}
A cereal box contains 1 of 4 different toys (equally likely). How many boxes must you buy to expect to collect all 4?
\end{problem}

\begin{solution}
The total number of boxes is a sum of four phase waits, so we use linearity of expectation. In phase $k$ we already hold $k-1$ toys and buy boxes until a new one appears. The per-trial success probability is $p = \frac{\text{favourable}}{\text{total}}$, the number of missing toys over 4. We want the reciprocal, the expected number of trials per success, which is $\frac{1}{p}$:
\begin{align*}
0 \text{ toys}, \; p = \tfrac{4}{4} \;&\longrightarrow\; \E[T_1] = \tfrac{1}{p} = 1 \\
1 \text{ toy}, \; p = \tfrac{3}{4} \;&\longrightarrow\; \E[T_2] = \tfrac{1}{p} = \tfrac{4}{3} \\
2 \text{ toys}, \; p = \tfrac{2}{4} \;&\longrightarrow\; \E[T_3] = \tfrac{1}{p} = 2 \\
3 \text{ toys}, \; p = \tfrac{1}{4} \;&\longrightarrow\; \E[T_4] = \tfrac{1}{p} = 4
\end{align*}
\[
\E[T] = \E[T_1] + \E[T_2] + \E[T_3] + \E[T_4] = 1 + \tfrac{4}{3} + 2 + 4 = \tfrac{25}{3} \approx 8.33
\]
\end{solution}

\begin{problem}
I roll a die repeatedly and keep a running sum. What's the expected number of rolls until the sum first reaches or exceeds 10?
\end{problem}

\begin{solution}
We solve recursively. Let $\E[s]$ be the expected number of additional rolls needed to first reach a sum of at least 10, starting from current sum $s$. The answer we want is $\E[0]$.

Base case: $\E[s] = 0$ for all $s \geq 10$, since the target is already met. For $s < 10$, one roll costs 1 and moves to $s+1, \dots, s+6$ with equal probability $\frac{1}{6}$:
\begin{align*}
\E[s] &= 1 + \frac{1}{6}\sum_{j=1}^{6} \E[s+j], \qquad \E[s] = 0 \text{ for } s \geq 10.
\end{align*}
Working downward from the base case:
\begin{align*}
\E[9] &= 1 + \tfrac{1}{6}\bigl(\E[10] + \cdots + \E[15]\bigr) = 1 \\
\E[8] &= 1 + \tfrac{1}{6}\E[9] = \tfrac{7}{6} \\
\E[7] &= 1 + \tfrac{1}{6}\bigl(\E[8] + \E[9]\bigr) = \tfrac{49}{36} \\
      &\;\;\vdots \\
\E[0] &\approx 3.6
\end{align*}
\end{solution}

\newpage

\begin{problem}
You roll a fair six-sided die. After seeing the result, you may keep it as your score, or reroll. If you reroll, the new value replaces the old one (you cannot go back to a discarded roll). You are allowed up to three rolls total, and your score is the value you are holding when you stop. Playing optimally to maximize
your expected score, what is that expected score, and what is the stopping rule at each stage?
\end{problem}

\begin{solution}
Let $\E[r]$ be the expected score with $r$ rerolls remaining. Solve backwards.

With no rerolls left, you keep whatever you roll, so $\E[0] = 3.5$.

With one reroll left, you keep the current face only if it beats the continuation value $\E[0] = 3.5$, i.e.\ keep on $\{4,5,6\}$ and reroll otherwise:
\[
\E[1] = \tfrac{1}{6}(4) + \tfrac{1}{6}(5) + \tfrac{1}{6}(6) + \tfrac{3}{6}(3.5) = \tfrac{17}{4} = 4.25.
\]
With two rerolls left, the continuation value is now $\E[1] = 4.25$, so you keep only faces that beat it, $\{5,6\}$, and reroll on $\{1,2,3,4\}$:
\[
\E[2] = \tfrac{1}{6}(6) + \tfrac{1}{6}(5) + \tfrac{4}{6}(4.25) = \tfrac{14}{3} \approx 4.67.
\]
\end{solution}


\begin{problem}
You pick a random number uniformly from [0, 1]. Then you pick another. Then another. You stop when the running total exceeds 1. How many numbers do you expect to pick?
\end{problem}

\begin{solution}
Let $N$ be the number of draws until the running sum first exceeds 1, and let $X_1, X_2, \dots$ be the draws. We want $\E[N]$.

\medskip
\textit{Step 1: a tail probability.} The \textit{tail probability} $\Prob(N > n)$ is the chance $N$ exceeds a cutoff $n$, i.e.\ the chance we still need more draws after $n$ of them. Needing more than $n$ draws means the first $n$ have not yet crossed 1:
\[
N > n \iff X_1 + \cdots + X_n \leq 1.
\]

\textit{Step 2: geometry} The point $(X_1, \dots, X_n)$ lands uniformly in the unit hypercube $[0,1]^n$, which has volume 1. Since the draws are uniform, the probability of landing in any region equals that region's volume. The region where the coordinates are nonnegative and sum to at most 1 is a \textit{simplex} (the $n$-dimensional triangle: a segment in 1D, triangle in 2D, tetrahedron in 3D). Its volume is $\frac{1}{n!}$ ($V_n(t) = \int_0^t V_{n-1}(t-x)\,dx$, starting from $V_1(t) = t$. Then $V_n(t) = \frac{t^n}{n!}$, evaluated at $t=1$, $V_n(1)= \frac{1}{n!}$), so
\[
\Prob(N > n) = \Prob(X_1 + \cdots + X_n \leq 1) = \frac{1}{n!}.
\]

\textit{Step 3: sum the tails.} For a nonnegative integer variable, the mean equals the sum of its tail probabilities:
\[
\E[N] = \sum_{n=0}^{\infty} \Prob(N > n).
\]
This holds because an outcome $N = k$ clears the cutoffs $n = 0, 1, \dots, k-1$, so it is counted exactly $k$ times across the sum, contributing its own value. Substituting the tails:
\[
\E[N] = \sum_{n=0}^{\infty} \frac{1}{n!} = \frac{1}{0!} + \frac{1}{1!} + \frac{1}{2!} + \cdots = e \approx 2.718.
\]
\end{solution}

\newpage

\subsection{Conditional Probability \& Bayes}

\begin{definition}[Conditional Probability]
The probability of $A$ given that $B$ has occurred is
\[
P(A \mid B) = \frac{P(A \cap B)}{P(B)}, \qquad P(B) > 0.
\]
Conditioning on $B$ restricts attention to the outcomes inside $B$ and rescales their probabilities to sum to 1, which is the role of the denominator.
\end{definition}

\begin{theorem}[Bayes' Theorem]
Bayes' theorem reverses the direction of conditioning, updating a hypothesis as evidence arrives:
\[
P(A \mid B) = \frac{P(B \mid A) \, P(A)}{P(B)}.
\]
The four pieces are worth memorizing by name:
\begin{itemize}[itemsep=-1ex, label=--]
  \item $P(A \mid B)$ is the posterior, the updated belief in $A$ after seeing $B$
  \item $P(B \mid A)$ is the likelihood, how well $A$ explains the evidence
  \item $P(A)$ is the prior, the belief in $A$ before the evidence
  \item $P(B)$ is the marginal, the total probability of the evidence
\end{itemize}
\end{theorem}

\begin{keybox}[Odds form skips the marginal]
Dividing Bayes' theorem for two competing hypotheses cancels the shared marginal $P(B)$:
\[
\underbrace{\frac{P(A \mid B)}{P(A^c \mid B)}}_{\text{posterior odds}}
= \underbrace{\frac{P(B \mid A)}{P(B \mid A^c)}}_{\text{likelihood ratio}}
\times \underbrace{\frac{P(A)}{P(A^c)}}_{\text{prior odds}}.
\]
Each piece of evidence multiplies the prior odds by its likelihood ratio, so repeated independent evidence just multiplies again. This avoids computing $P(B)$ and is the fastest route through two-hypothesis update problems.
\end{keybox}

\medskip

\begin{theorem}[Law of Total Probability]
If $B_1, B_2, \dots , B_n$ partition the sample space (mutually exclusive and exhaustive), then
\[
P(A) = \sum_i P(A \mid B_i) \, P(B_i).
\]
\end{theorem}

This is how the marginal in Bayes' theorem is usually assembled. When $P(A)$ is not known directly but the conditional probabilities on each case are, summing those conditionals weighted by the case probabilities builds it. The two-way form $P(A) = P(A \mid C)P(C) + P(A \mid C^c)P(C^c)$ covers most interview uses.

\medskip

\begin{example}[Monty Hall Problem]
Assume the contestant chooses Door 1 and Monty opens Door 2 to reveal a goat. Let $A_i$ be the event the car is behind Door $i$ (prior $P(A_i) = \frac{1}{3}$), and $B$ be the event Monty opens Door 2.

The likelihoods $P(B|A_i)$ are:
\begin{itemize}[itemsep=-1ex, label=--]
    \item $P(B|A_1) = \frac{1}{2}$ (Monty randomly picks Door 2 or 3)
    \item $P(B|A_2) = 0$ (Monty cannot reveal the car)
    \item $P(B|A_3) = 1$ (Monty is forced to open Door 2)
\end{itemize}

The marginal probability of $B$ is:
\[
P(B) = \sum_{i=1}^3 P(B|A_i)P(A_i) = \left(\frac{1}{2}\right)\left(\frac{1}{3}\right) + (0)\left(\frac{1}{3}\right) + (1)\left(\frac{1}{3}\right) = \frac{1}{2}
\]

By Bayes' Theorem, the posterior probabilities of winning by staying ($A_1$) or switching ($A_3$) are:
\[
P(A_1|B) = \frac{P(B|A_1)P(A_1)}{P(B)} = \frac{\left(\frac{1}{2}\right)\left(\frac{1}{3}\right)}{\frac{1}{2}} = \frac{1/6}{1/2} = \frac{1}{3}
\]
\[
P(A_3|B) = \frac{P(B|A_3)P(A_3)}{P(B)} = \frac{(1)\left(\frac{1}{3}\right)}{\frac{1}{2}} = \frac{1/3}{1/2} = \frac{2}{3}
\]

Switching doubles the probability of winning.
\end{example}

\newpage

\subsubsection{Select Worked Problems}

\begin{problem}
I roll two dice. Given that the sum is 8, what's the probability that one of the dice shows a 3?
\end{problem}

\begin{solution}
The possible ways to roll a sum of 8 are: $(2, 6), (3, 5), (4, 4), (5, 3), (6, 2)$. Of these, two options have a 3 in them, so the answer is $\frac{2}{5} = 0.4$.
\end{solution}

\begin{problem}
A bag has 3 red and 2 blue balls. You draw one, don't replace it, then draw another. What's P(2nd is red | 1st was blue)?
\end{problem}

\begin{solution}
Once we drew one blue on, the remaining bag composition is 3 red balls and 1 blue one. Then, there is a $\frac{3}{4}$ probability that the second will be red.
\end{solution}


\begin{problem}
60\% of emails are spam. A filter catches 95\% of spam but also flags 3\% of legitimate emails. An email is flagged. What's the probability it's actually spam?
\end{problem}

\begin{solution}
We want $P(Spam | Flagged)$. The likelihood $P(Flagged | Spam)$ is given as $0.95$, the prior $P(Spam)$ is $0.6$, and the marginal probability $P(Flagged)$ is $0.95 \cdot 0.6 + 0.03 \cdot 0.4 = 0.582$.

Then the posterior is
\[
P(Spam | Flagged) = \frac{P(Flagged | Spam) \times P(Spam)}{P(Flagged)} = \frac{0.95 \times 0.6}{0.582} \approx 0.979
\]


\end{solution}


\begin{problem}
You have 2 coins. Coin A is fair (50\% heads). Coin B is biased (75\% heads). You pick a coin at random and flip it. It lands heads. What's the probability you picked coin B?
\end{problem}

\begin{solution}
We are looking for the posterior, $P(B|H)$.

The priors are $P(A) = P(B) = 0.5$. The likelihoods are $P(H|A) = 0.5$ and $P(H|B) = 0.75$. Then $P(H) = P(H|A) \cdot P(A) + P(H|B) \cdot P(B) = 0.5 \cdot 0.5 + 0.75 \cdot 0.5 = 0.625$. Then finally
\[
P(B|H) = \frac{P(H|B) \cdot P(B)}{P(H)} = \frac{0.75 \cdot 0.5}{0.625} = 0.6
\]
After flipping and landing on heads, our belief that we picked coin B increased from $0.5$ to $0.6$.
\end{solution}

\begin{problem}
Same as 1.31. The chosen coin is flipped twice and get heads both times. Now what's $P(B)$?
\end{problem}

\begin{solution}
The only difference is our priors are $P(A) = 0.4$ and $P(B) = 0.6$. Then we update $P(H) = P(H|A) \cdot P(A) + P(H|B) \cdot P(B) = 0.5 \cdot 0.4 + 0.75 \cdot 0.6 = 0.65$. Finally:
\[
P(B|H) = \frac{P(H|B) \cdot P(B)}{P(H)} = \frac{0.75 \cdot 0.6}{0.65} = \frac{9}{13} \approx 0.692
\]
\end{solution}

\newpage

\begin{problem}
A family has 2 children. Given that at least one child is a boy, what's the probability both are boys?
\end{problem}

\begin{solution}
We can answer this by considering the restriction over the sample space caused by the new information. Originally, for two children there would be four possibilities, of which 1 has two boys. After we know that at least one child is a boy, only three options remain $(B,G), (B,B), (G, B)$. Thus, the probability both are boys is $\frac{1}{3}$.
\end{solution}

\begin{problem}
Three cards: one is red on both sides, one is blue on both sides, one is red on one side and blue on the other. You pick a card at random and see a red side. What's the probability the other side is also red?
\end{problem}

\begin{solution}
The natural first instinct is to treat the cards as the sample space, as with the births problem: restrict to cards consistent with seeing a red face, then count outcomes among those. This fails here because a card is not a single unit of observation. When a card has two sides of the same color, ``seeing red'' can happen in more than one way for that card. There's no such internal multiplicity in the births problem.

So take the six sides, not the three cards, as the equally likely outcomes. Three of these six are red: two belong to the $RR$ card, one belongs to the $RB$ card. Since each red side was equally likely to be the one we observed, and two of the three red sides have a red back:
\[
P(\text{other side red} \mid \text{observed red}) = \frac{2}{3}.
\]
\end{solution}


\begin{problem}
You have a 6-sided die. You don't know if it's fair or loaded. A fair die shows each number with probability 1/6. The loaded die shows 6 with probability 1/2 and each other number with probability 1/10. You pick a die at random (50-50) and roll a 6. What's P(loaded)?
\end{problem}

\begin{solution}
By Bayes' theorem, with $P(\text{loaded}) = P(\text{fair}) = 0.5$:
\[
P(\text{loaded} \mid 6) = \frac{P(6 \mid \text{loaded}) \, P(\text{loaded})}{P(6)} = \frac{0.5 \times 0.5}{0.5 \times 0.5 + \tfrac{1}{6} \times 0.5} = \frac{0.25}{1/3} = 0.75
\]

Alternatively, work in odds instead of probabilities. Dividing the Bayes' theorem expressions for $P(\text{loaded} \mid 6)$ and $P(\text{fair} \mid 6)$ cancels $P(6)$ from both, leaving:
\[
\underbrace{\frac{P(\text{loaded} \mid 6)}{P(\text{fair} \mid 6)}}_{\text{posterior odds}} = \underbrace{\frac{P(6 \mid \text{loaded})}{P(6 \mid \text{fair})}}_{\text{likelihood ratio}} \times \underbrace{\frac{P(\text{loaded})}{P(\text{fair})}}_{\text{prior odds}}
\]
So the posterior odds are just the prior odds scaled by the likelihood ratio, no need to compute $P(6)$ at all. Here, the prior odds are $\frac{P(\text{loaded})}{P(\text{fair})} = \frac{0.5}{0.5} = 1$, and the likelihood ratio is $\frac{P(6 \mid \text{loaded})}{P(6 \mid \text{fair})} = \frac{0.5}{1/6} = 3$. So the posterior odds are $1 \times 3 = 3$, i.e. $3:1$ in favor of loaded, giving $P(\text{loaded} \mid 6) = 3/4 = 0.75$.

In short, a six is three times more likely to appear if the die is loaded, so each 6 observed multiplies the odds by 3. If we observed \textit{another} 6, the odds would update again: $3:1 \rightarrow 9:1$, giving $P(\text{loaded}) = \frac{9}{10}$.

\end{solution}

\newpage

% ============================================================
\subsection{Distributions}
% ============================================================

A named distribution is a compressed description of a random process. What makes the catalogue useful is not the formulas in isolation but the mapping from a verbal cue in the problem statement to a process, and from that process to a mean, a variance, and a tail. The organizing question for any problem is therefore which process generated the randomness rather than which formula to apply.

\begin{definition}[Random variable and distribution]
A random variable $X$ is a function assigning a number to each outcome in the sample space $\Omega$. Its distribution specifies how probability is spread across the values $X$ can take, and it is described by one of three equivalent objects:
\begin{itemize}[itemsep=-1ex, label=--]
  \item the probability mass function (PMF) $P(X = x)$ in the discrete case, with $\sum_x P(X=x) = 1$
  \item the probability density function (PDF) $f(x)$ in the continuous case, with $\int_{-\infty}^{\infty} f(x)\,dx = 1$
  \item the cumulative distribution function (CDF) $F(x) = P(X \leq x)$, defined in both cases
\end{itemize}
The set of values with nonzero probability is the support.
\end{definition}

In the continuous case $P(X = x) = 0$ for every single point, so density is not probability. Only integrals of the density over an interval carry probability, which is why continuous questions are answered through the CDF as $P(a < X \leq b) = F(b) - F(a)$.

\vspace{1em}

\begin{keybox}[Recognize the process first]
A named distribution (Bernoulli, Binomial, Geometric, Poisson, and so on) is a family arising from a common underlying random process, indexed by one or more parameters such as $n$ and $p$ for the Binomial. Once a problem is matched to its process, the PMF or PDF, mean, and variance follow immediately from the known formulas for the family. Three cues carry most of the discrete cases:
\begin{itemize}[itemsep=0.5ex, label=--]
  \item a \textit{fixed} number of trials with a counted number of successes points to the Binomial
  \item a \textit{fixed} number of successes with a counted number of trials points to the Geometric or Negative Binomial
  \item a count over a window of time or space with no natural trial count points to the Poisson
\end{itemize}
\end{keybox}

\begin{definition}[Moment]
The $n$-th moment of a random variable $X$ is $\E[X^n]$, and its $n$-th central moment is $\E[(X - \mu)^n]$, taken about the mean $\mu = \E[X]$. A moment exists only when the defining expectation converges absolutely, meaning $\E[|X|^n] < \infty$.
\end{definition}

Moments are summary numbers describing the shape of a distribution, each one weighting outcomes by a higher power and so responding more strongly to values far from the center. The first four carry standard names and interpretations:
\begin{itemize}[itemsep=-1ex, label=--]
  \item the first moment $\E[X]$ is the mean, locating the distribution
  \item the second central moment $\E[(X-\mu)^2]$ is the variance, measuring spread
  \item the third standardized moment is the skewness, measuring asymmetry
  \item the fourth standardized moment is the kurtosis, measuring tail weight
\end{itemize}
Standardizing divides the central moment by $\sigma^n$, which makes skewness and kurtosis dimensionless and therefore comparable across distributions. A distribution is not guaranteed to possess any given moment, since sufficiently heavy tails make the integral diverge, and moments fail from the top down in the sense that a finite $n$-th moment forces every lower moment to be finite as well.

% ============================================================
\subsubsection{Bernoulli Distribution}
% ============================================================

\begin{definition}[Bernoulli]
A random variable $X$ follows a Bernoulli distribution with parameter $p \in [0,1]$, written $X \sim \text{Bernoulli}(p)$, if it takes exactly two values:
\[
P(X = 1) = p, \qquad P(X = 0) = 1 - p.
\]
It models a single trial with two outcomes, conventionally labeled success ($1$) and failure ($0$).
\[
\mu = p \hspace{5em} \sigma^2 = p(1-p)
\]
\end{definition}

Both moments follow in one line. Since $X$ only takes the values $0$ and $1$, the identity $X^2 = X$ holds pointwise, so $\E[X^2] = \E[X] = p$ and
\[
\Var(X) = \E[X^2] - \E[X]^2 = p - p^2 = p(1-p).
\]
The variance is largest at $p = \tfrac{1}{2}$, where the outcome is most uncertain, and shrinks to zero as $p$ approaches $0$ or $1$, where the result becomes nearly certain.

\vspace{1em}

\begin{keybox}[Building block for other distributions]
The Bernoulli is rarely the distribution of interest on its own. It is the atomic unit that other named distributions are built from. A Binomial random variable is a sum of $n$ i.i.d.\ Bernoulli$(p)$ trials, and a Geometric random variable counts trials until the first Bernoulli success. Recognizing a Bernoulli inside a larger problem, often through an indicator variable $\ind\{\cdot\}$, is frequently the first step in applying linearity of expectation to more complex sums.
\end{keybox}

% ============================================================
\subsubsection{Binomial Distribution}
% ============================================================

\begin{definition}[Binomial]
$X \sim \text{Binomial}(n,p)$ counts the number of successes in $n$ independent trials, each succeeding with the same probability $p$. Its PMF is
\[
P(X=k) = \binom{n}{k} p^k (1-p)^{n-k}, \qquad k = 0, 1, \dots, n.
\]
\end{definition}

\begin{proof}
Fix one particular sequence of outcomes with $k$ successes and $n-k$ failures. Independence makes its probability the product $p^k (1-p)^{n-k}$, and that value depends only on the counts, not on the positions. The number of distinct sequences with $k$ successes is the number of ways to choose which trials succeed, namely $\binom{n}{k}$. Summing the common probability over those disjoint sequences gives the stated PMF.
\end{proof}
\[
\mu = np \hspace{5em} \sigma^2 = np(1-p)
\]

Writing $X = \sum_{i=1}^n X_i$ as a sum of i.i.d.\ Bernoulli$(p)$ variables gives both moments without touching the PMF. Linearity of expectation yields $\E[X] = np$ with no independence needed. Variance requires independence, which zeroes every covariance cross term and leaves $\Var(X) = n \cdot p(1-p)$.

\newpage

\begin{keybox}[Sampling with replacement against without]
The Binomial requires that $p$ stay constant across trials, which holds when sampling with replacement. Sampling without replacement changes the composition of the pool after each draw and produces the Hypergeometric distribution instead. When the population is large relative to the sample, the shift in $p$ is negligible and the Binomial is the standard approximation.
\end{keybox}

\begin{definition}[Hypergeometric]
Drawing $n$ items without replacement from a population of $N$ containing $K$ successes, the number of successes drawn has PMF
\[
P(X = k) = \frac{\binom{K}{k}\binom{N-K}{n-k}}{\binom{N}{n}}.
\]
Its mean is $\mu = n\frac{K}{N}$, matching the Binomial with $p = K/N$, and its variance is
\[
\sigma^2 = np(1-p)\,\frac{N-n}{N-1},
\]
where the factor $\frac{N-n}{N-1}$ is the finite population correction.
\end{definition}

The mean is unchanged because linearity of expectation ignores the dependence between draws. The variance is lower because the draws are negatively correlated: each draw removes an item from the pool, so a success on one draw leaves fewer successes available and makes a success on the next slightly less likely. That downward pull is exactly the correction factor, which is below 1, so sampling without replacement is strictly less variable than sampling with replacement. It collapses to 1 as $N$ grows relative to $n$, which recovers the Binomial.

% ============================================================
\subsubsection{Geometric Distribution}
% ============================================================

\begin{definition}[Geometric]
$X \sim \text{Geometric}(p)$ counts the number of independent trials up to and including the first success, where each trial succeeds with probability $p$. Its PMF is
\[
P(X = k) = (1-p)^{k-1} p, \qquad k = 1, 2, 3, \dots
\]
\end{definition}

The PMF is a direct product. Reaching the first success on trial $k$ requires $k-1$ failures followed by one success, and independence multiplies those probabilities. The tail probability is equally compact, since needing more than $k$ trials means the first $k$ all failed:
\[
P(X > k) = (1-p)^k.
\]

\[
\mu = \frac{1}{p} \hspace{5em} \sigma^2 = \frac{1-p}{p^2}
\]

\begin{keybox}[Two conventions]
Some texts define the Geometric as the number of \textit{failures before} the first success, with support $k = 0, 1, 2, \dots$ and PMF $(1-p)^k p$. That version has mean $\frac{1-p}{p}$, exactly one less, while the variance $\frac{1-p}{p^2}$ is identical since shifting a variable by a constant does not change its spread. The two differ by the single trial on which the success occurs, so the convention in force must be established before any numerical comparison. The trials convention with mean $\frac{1}{p}$ is the one used throughout these notes.
\end{keybox}

\begin{theorem}[Memorylessness]
The Geometric is the only discrete distribution with the memoryless property:
\[
P(X > m + n \mid X > m) = P(X > n).
\]
Past failures carry no information about how many further trials are required.
\end{theorem}

\begin{proof}
Apply the definition of conditional probability. The event $\{X > m+n\}$ is contained in $\{X > m\}$, so their intersection is the former:
\[
P(X > m+n \mid X > m) = \frac{P(X > m+n)}{P(X > m)} = \frac{(1-p)^{m+n}}{(1-p)^m} = (1-p)^n = P(X > n). \qedhere
\]
\end{proof}

Memorylessness is the reason the recursive expectation argument works. A coin that has landed tails ten times in a row still needs $\frac{1}{p}$ further flips in expectation, so conditioning on the first trial reproduces the original problem and gives $\E[X] = 1 + (1-p)\E[X]$, which solves to $\frac{1}{p}$.

% ============================================================
\subsubsection{Negative Binomial Distribution}
% ============================================================

\begin{definition}[Negative Binomial]
$X \sim \text{NegBin}(r,p)$ counts the number of independent trials up to and including the $r$-th success. Its PMF is
\[
P(X = k) = \binom{k-1}{r-1} p^r (1-p)^{k-r}, \qquad k = r, r+1, \dots
\]
\end{definition}

The binomial coefficient counts the arrangements of the first $r-1$ successes among the first $k-1$ trials. Trial $k$ is forced to be a success, which is what pins the $r$-th success to the final position rather than anywhere in the sequence.

\[
\mu = \frac{r}{p} \hspace{5em} \sigma^2 = \frac{r(1-p)}{p^2}
\]

Both follow from writing $X = G_1 + G_2 + \cdots + G_r$, a sum of $r$ independent Geometric$(p)$ waiting times. Memorylessness is what guarantees each waiting period restarts fresh after a success, so the $G_i$ are i.i.d. The mean is $r$ times the Geometric mean by linearity of expectation. The variance is $r$ times the Geometric variance by \textit{independence}, not by linearity, since variances add only when covariance terms vanish.

% ============================================================
\subsubsection{Poisson Distribution}
% ============================================================

\begin{definition}[Poisson]
$X \sim \text{Poisson}(\lambda)$ counts discrete events occurring in a fixed window of time or space, when events arrive independently at a constant average rate $\lambda$ per window. Its PMF is
\[
P(X = k) = e^{-\lambda}\frac{\lambda^k}{k!}, \qquad k = 0, 1, 2, \dots
\]
\end{definition}

\[
\mu = \lambda \hspace{5em} \sigma^2 = \lambda
\]

The equality of mean and variance serves as a diagnostic. Count data whose sample variance materially exceeds its mean is overdispersed and is not Poisson. This occurs routinely in market data, where order arrivals cluster rather than arriving at a constant rate, and it motivates the mixture and doubly stochastic models that generalize the Poisson.

\begin{theorem}[Poisson limit of the Binomial]
Fix $\lambda > 0$ and let $n \to \infty$ with $p = \lambda/n$, so that the expected count $np = \lambda$ stays constant. Then
\[
\binom{n}{k} p^k (1-p)^{n-k} \longrightarrow e^{-\lambda}\frac{\lambda^k}{k!}.
\]
\end{theorem}

\begin{proof}
Substitute $p = \lambda/n$ and split the expression into three factors:
\[
\underbrace{\frac{n(n-1)\cdots(n-k+1)}{n^k}}_{\to\, 1}
\cdot \frac{\lambda^k}{k!}
\cdot \underbrace{\left(1 - \frac{\lambda}{n}\right)^{n}}_{\to\, e^{-\lambda}}
\cdot \underbrace{\left(1 - \frac{\lambda}{n}\right)^{-k}}_{\to\, 1}.
\]
The first factor is a ratio of $k$ terms each approaching $n$, the third is the standard limit defining $e^{-\lambda}$, and the fourth has a fixed exponent with a base approaching 1. The surviving product is the Poisson PMF.
\end{proof}

The interpretation matters more than the algebra. The Poisson is the law of rare events, meaning many opportunities for an event to occur with a small probability at each one. This is why it models message counts, quote updates, and defaults, where no natural trial count exists but a rate does.

\vspace{1ex}

\begin{keybox}[Superposition and thinning]
\textbf{Superposition.} Independent Poisson streams merge into a Poisson stream with the summed rate, so $X \sim \text{Poisson}(\lambda_1)$ and $Y \sim \text{Poisson}(\lambda_2)$ independent give $X + Y \sim \text{Poisson}(\lambda_1 + \lambda_2)$. Combining order flow across venues is the standard application.

\vspace{2ex}

\textbf{Thinning.} If each event in a Poisson$(\lambda)$ stream is independently retained with probability $q$, the retained events form a Poisson$(\lambda q)$ stream, and the discarded ones form an independent Poisson$(\lambda(1-q))$ stream. Filtering a stream to buy orders only is the standard application.
\end{keybox}

% ============================================================
\subsubsection{Exponential Distribution}
% ============================================================

\begin{definition}[Exponential]
$X \sim \text{Exponential}(\lambda)$ is the continuous waiting time until the next event in a Poisson process of rate $\lambda$. Its density and CDF are
\[
f(x) = \lambda e^{-\lambda x}, \qquad F(x) = 1 - e^{-\lambda x}, \qquad x \geq 0.
\]
\end{definition}

\[
\mu = \frac{1}{\lambda} \hspace{5em} \sigma^2 = \frac{1}{\lambda^2}
\]

The Exponential is the continuous counterpart of the Geometric, and the pairing runs deep. The Geometric counts discrete trials until a success while the Exponential measures continuous time until an arrival, and each is the unique memoryless distribution on its domain:
\[
P(X > s + t \mid X > s) = P(X > t).
\]
The proof is the same cancellation as before, using the tail $P(X > x) = e^{-\lambda x}$. The consequence is that the expected remaining wait is $\frac{1}{\lambda}$ regardless of how long the wait has already lasted, so elapsed time carries no information about the time still to come.

\begin{keybox}[Counts against waiting times]
Poisson and Exponential describe one process from two directions. If events arrive at rate $\lambda$, the \textit{count} in a unit window is Poisson$(\lambda)$ and the \textit{gap} between consecutive events is Exponential$(\lambda)$. The identity linking them is that zero arrivals in $[0,t]$ is the same event as a wait longer than $t$:
\[
P(N_t = 0) = e^{-\lambda t} = P(X > t).
\]
\textbf{Minimum of exponentials.} If $X_i \sim \text{Exponential}(\lambda_i)$ independently, then $\min_i X_i \sim \text{Exponential}\!\left(\sum_i \lambda_i\right)$, and the probability that stream $j$ arrives first is $\lambda_j / \sum_i \lambda_i$. This is the central result for competing arrivals, such as determining which of several order flows reaches the book first.
\end{keybox}

% ============================================================
\subsubsection{Uniform Distribution}
% ============================================================

\begin{definition}[Continuous Uniform]
$X \sim \text{Uniform}(a,b)$ places equal density across an interval:
\[
f(x) = \frac{1}{b-a}, \qquad F(x) = \frac{x-a}{b-a}, \qquad a \leq x \leq b.
\]
\end{definition}

\[
\mu = \frac{a+b}{2} \hspace{5em} \sigma^2 = \frac{(b-a)^2}{12}
\]

Probability for a continuous Uniform is proportional to length, so $P(c < X < d) = \frac{d-c}{b-a}$ for any subinterval. The variance depends only on the width $b-a$ and not on the location, which is the general behavior of variance under a shift.

\begin{keybox}[Do not mix the discrete and continuous forms]
The discrete Uniform on $\{1, 2, \dots, n\}$ has mean $\frac{n+1}{2}$ and variance $\frac{n^2 - 1}{12}$, which is not the continuous formula with $a=1$ and $b=n$. A fair die has $\mu = 3.5$ and $\sigma^2 = \frac{35}{12} \approx 2.92$. Substituting into $\frac{(b-a)^2}{12}$ would give $\frac{25}{12}$, which is wrong.
\end{keybox}

\begin{keybox}[Universality of the uniform]
For any continuous random variable $X$ with CDF $F$, the transformed variable $F(X)$ is Uniform$(0,1)$. Running this backwards, $F^{-1}(U)$ with $U \sim \text{Uniform}(0,1)$ has distribution $F$. This inverse transform is the basis for generating arbitrary distributions from a uniform random number generator, which is the foundation of Monte Carlo simulation. For the Exponential it gives the explicit sampler $X = -\frac{1}{\lambda}\ln(U)$.
\end{keybox}

% ============================================================
\subsubsection{Normal (Gaussian) Distribution}
% ============================================================

\begin{definition}[Normal]
$X \sim \mathcal{N}(\mu, \sigma^2)$ has density
\[
f(x) = \frac{1}{\sqrt{2\pi\sigma^2}}\exp\!\left(-\frac{(x-\mu)^2}{2\sigma^2}\right).
\]
The two parameters are exactly its mean and its variance, so the distribution is fully specified by its first two moments.
\end{definition}

Two structural properties account for most applications.

\textbf{Standardization.} Any Normal reduces to the standard Normal through $Z = \frac{X - \mu}{\sigma} \sim \mathcal{N}(0,1)$, which is why a single table of $Z$ values suffices for every Normal question.

\textbf{Closure under addition.} A sum of independent Normals is Normal, with means and variances adding:
\[
X_1 \sim \mathcal{N}(\mu_1, \sigma_1^2), \; X_2 \sim \mathcal{N}(\mu_2, \sigma_2^2) \implies X_1 + X_2 \sim \mathcal{N}(\mu_1 + \mu_2, \, \sigma_1^2 + \sigma_2^2).
\]
Standard deviations do not add. This is the source of the $\sqrt{T}$ scaling of volatility, since variance grows linearly in time while volatility grows with the square root.

\vspace{1em}

\begin{keybox}[Standard tail masses]
\begin{center}
\begin{tabular}{lll}
\toprule
\textbf{Range} & \textbf{Mass inside} & \textbf{Mass outside} \\
\midrule
$\mu \pm 1\sigma$      & $68.3\%$ & $\approx 1$ in $3$ \\
$\mu \pm 1.645\sigma$  & $90\%$   & $5\%$ per tail \\
$\mu \pm 1.96\sigma$   & $95\%$   & $2.5\%$ per tail \\
$\mu \pm 2\sigma$      & $95.4\%$ & $\approx 1$ in $22$ \\
$\mu \pm 2.576\sigma$  & $99\%$   & $0.5\%$ per tail \\
$\mu \pm 3\sigma$      & $99.7\%$ & $\approx 1$ in $370$ \\
\bottomrule
\end{tabular}
\end{center}

\vspace{2ex}

The one-tailed values $1.645$ and $2.326$ correspond to the $95$th and $99$th percentiles and are the multipliers underlying standard value-at-risk calculations.
\end{keybox}

\vspace{1ex}

\begin{remark}[Why the Normal is the default]
The Normal is the maximum entropy distribution among all distributions on the real line with a given mean and variance. Constraining only the first two moments and otherwise assuming as little as possible produces the Normal, which is the formal version of the claim that it encodes minimal structural assumptions. The Central Limit Theorem gives the second, more operational reason for its ubiquity.
\end{remark}

\begin{theorem}[Central Limit Theorem]
Let $X_1, \dots, X_n$ be i.i.d.\ with finite mean $\mu$ and finite variance $\sigma^2$. Then the standardized sample mean converges in distribution to a standard Normal:
\[
\frac{\bar{X}_n - \mu}{\sigma/\sqrt{n}} \xrightarrow{\;d\;} \mathcal{N}(0,1)
\]
\end{theorem}

The standardization is essential to the statement. The sample mean $\bar{X}_n$ on its own converges to the constant $\mu$ by the law of large numbers, and its variance $\sigma^2/n$ shrinks to zero, so it is the rescaled deviation that stabilizes into a Normal shape rather than the mean itself.

\vspace{1ex}

\begin{keybox}[Where the CLT fails]
The common rule that $n \geq 30$ suffices is a heuristic, not a theorem, and it breaks in three ways that matter for trading:
\begin{itemize}[itemsep=1ex, label=--]
  \item \textbf{Finite variance is required.} Heavy-tailed distributions with infinite variance, such as the Cauchy, never converge. The sample mean of Cauchy draws has the same Cauchy distribution as a single draw, so averaging buys nothing.
  \item \textbf{Skewness slows convergence.} Strongly asymmetric or rare-event distributions may need $n$ in the thousands before the Normal approximation is usable in the tails, which is exactly where risk questions live.
  \item \textbf{Independence is required.} Financial returns are volatility clustered and serially dependent, so effective sample size is far below nominal sample size.
\end{itemize}

\vspace{1ex}

Empirical asset returns are leptokurtic, meaning excess kurtosis above the Normal value of 3, so extreme moves occur far more often than a Gaussian model predicts.
\end{keybox}

% ============================================================
\subsubsection{Lognormal Distribution}
% ============================================================

\begin{definition}[Lognormal]
$X$ is lognormal when $\ln X \sim \mathcal{N}(\mu, \sigma^2)$, so $X = e^{Y}$ for a Normal $Y$. Its support is $x > 0$, and
\[
\E[X] = e^{\mu + \sigma^2/2}, \qquad \text{median}(X) = e^{\mu}, \qquad \Var(X) = \left(e^{\sigma^2} - 1\right)e^{2\mu + \sigma^2}.
\]
\end{definition}

The lognormal is the standard model for asset prices, for two structural reasons. Prices are bounded below by zero while a Normal would assign positive probability to negative prices, and returns compound multiplicatively, so the logarithm of the price is a sum of log returns and therefore Normal by the same CLT argument.

The mean sits above the median because the distribution is right-skewed. Exponentiating a Normal stretches the upper values far more than the lower ones, producing a long right tail. The median only tracks the middle draw, which is $e^{\mu}$, but the mean is dragged upward by that tail to $e^{\mu + \sigma^2/2}$. The two differ by the factor $e^{\sigma^2/2}$, which sits at 1 when there is no variability and grows as the volatility $\sigma$ increases. More spread in the log means a heavier right tail and a wider gap between the average outcome and the typical one.

This gap is a concrete case of a general rule: averaging after applying a curved function is not the same as applying the function to the average. Here the function is $e^x$, which curves upward, so $\E[e^Y]$ comes out larger than $e^{\E[Y]}$. The excess is exactly the $e^{\sigma^2/2}$ term. The same correction reappears in the drift of geometric Brownian motion, the standard continuous-time price model, for the same reason.

% ============================================================
\subsubsection{Reference Tables}
% ============================================================

\begin{keybox}[Discrete distributions]
\begin{center}
\footnotesize
\setlength{\tabcolsep}{3pt}
\begin{tabular}{@{}lllccl@{}}
\toprule
\textbf{Distribution} & \textbf{Cue in the problem} & \textbf{PMF} & \textbf{Mean} & \textbf{Variance} & \textbf{Support} \\
\midrule
Bernoulli$(p)$ & one trial, success or failure & $p^x(1-p)^{1-x}$ & $p$ & $p(1-p)$ & $\{0,1\}$ \\[0.35em]
Binomial$(n,p)$ & successes in $n$ fixed trials & $\binom{n}{k}p^k(1-p)^{n-k}$ & $np$ & $np(1-p)$ & $0,\dots,n$ \\[0.35em]
Hypergeom. & successes, no replacement & $\dfrac{\binom{K}{k}\binom{N-K}{n-k}}{\binom{N}{n}}$ & $np$ & $np(1-p)\frac{N-n}{N-1}$ & $0,\dots,n$ \\[0.5em]
Geometric$(p)$ & trials until 1st success & $(1-p)^{k-1}p$ & $\dfrac{1}{p}$ & $\dfrac{1-p}{p^2}$ & $1,2,\dots$ \\[0.5em]
NegBin$(r,p)$ & trials until $r$-th success & $\binom{k-1}{r-1}p^r(1-p)^{k-r}$ & $\dfrac{r}{p}$ & $\dfrac{r(1-p)}{p^2}$ & $r, r+1,\dots$ \\[0.5em]
Poisson$(\lambda)$ & count at a rate, no trial count & $e^{-\lambda}\dfrac{\lambda^k}{k!}$ & $\lambda$ & $\lambda$ & $0,1,2,\dots$ \\[0.35em]
Uniform$\{1,n\}$ & equally likely integers & $\dfrac{1}{n}$ & $\dfrac{n+1}{2}$ & $\dfrac{n^2-1}{12}$ & $1,\dots,n$ \\
\bottomrule
\end{tabular}
\end{center}
\end{keybox}

\vspace{2ex}

\begin{keybox}[Continuous distributions]
\begin{center}
\footnotesize
\setlength{\tabcolsep}{3pt}
\begin{tabular}{@{}lllccl@{}}
\toprule
\textbf{Distribution} & \textbf{Cue in the problem} & \textbf{PDF} & \textbf{Mean} & \textbf{Variance} & \textbf{Support} \\
\midrule
Uniform$(a,b)$ & equally likely over a range & $\dfrac{1}{b-a}$ & $\dfrac{a+b}{2}$ & $\dfrac{(b-a)^2}{12}$ & $[a,b]$ \\[0.5em]
Exponential$(\lambda)$ & wait until next arrival & $\lambda e^{-\lambda x}$ & $\dfrac{1}{\lambda}$ & $\dfrac{1}{\lambda^2}$ & $x \geq 0$ \\[0.5em]
Normal$(\mu,\sigma^2)$ & sum of many small effects & $\frac{1}{\sqrt{2\pi\sigma^2}}e^{-(x-\mu)^2/2\sigma^2}$ & $\mu$ & $\sigma^2$ & $\mathbb{R}$ \\[0.5em]
Lognormal$(\mu,\sigma^2)$ & multiplicative growth & $\frac{1}{x\sqrt{2\pi\sigma^2}}e^{-(\ln x-\mu)^2/2\sigma^2}$ & $e^{\mu+\sigma^2/2}$ & $(e^{\sigma^2}\!-\!1)e^{2\mu+\sigma^2}$ & $x > 0$ \\
\bottomrule
\end{tabular}
\end{center}
\end{keybox}

\begin{keybox}[Moment generating functions]
The MGF is $M_X(t) = \E[e^{tX}]$. It determines the distribution uniquely, generates moments through $\E[X^n] = M_X^{(n)}(0)$, and multiplies across independent sums, which is the fastest proof that Binomials and Poissons and Normals are each closed under addition.
\begin{center}
\footnotesize
\begin{tabular}{@{}ll@{\hskip 3em}ll@{}}
\toprule
\textbf{Distribution} & $M_X(t)$ & \textbf{Distribution} & $M_X(t)$ \\
\midrule
Bernoulli$(p)$ & $1-p+pe^t$ & Poisson$(\lambda)$ & $e^{\lambda(e^t-1)}$ \\[0.4em]
Binomial$(n,p)$ & $(1-p+pe^t)^n$ & Exponential$(\lambda)$ & $\dfrac{\lambda}{\lambda-t}, \; t<\lambda$ \\[0.5em]
Geometric$(p)$ & $\dfrac{pe^t}{1-(1-p)e^t}$ & Normal$(\mu,\sigma^2)$ & $e^{\mu t + \sigma^2 t^2/2}$ \\[0.5em]
NegBin$(r,p)$ & $\left(\dfrac{pe^t}{1-(1-p)e^t}\right)^{\!r}$ & Lognormal & undefined for $t>0$ \\
\bottomrule
\end{tabular}
\end{center}
The lognormal entry is not an oversight. Its right tail is heavy enough that $\E[e^{tX}]$ diverges for every $t > 0$, which is a compact example of a distribution that has all moments yet no MGF.
\end{keybox}

\vspace{2ex}

\begin{keybox}[How the families connect]
\begin{center}
\footnotesize
\begin{tabular}{@{}lll@{}}
\toprule
\textbf{From} & \textbf{Relationship} & \textbf{To} \\
\midrule
Bernoulli & sum of $n$ i.i.d.\ trials & Binomial \\[0.3em]
Bernoulli & count trials to 1st success & Geometric \\[0.3em]
Geometric & sum of $r$ i.i.d.\ waits & Negative Binomial \\[0.3em]
Binomial & $n \to \infty$, $p \to 0$, $np = \lambda$ fixed & Poisson \\[0.3em]
Binomial & $n$ large, $np(1-p)$ large & Normal \\[0.3em]
Poisson & gap between arrivals & Exponential \\[0.3em]
Geometric & continuous time limit & Exponential \\[0.3em]
Exponential & minimum of independent draws & Exponential (rates add) \\[0.3em]
Normal & exponentiate & Lognormal \\[0.3em]
Any distribution & standardized mean, finite variance & Normal (CLT) \\[0.3em]
Any continuous & apply own CDF & Uniform$(0,1)$ \\
\bottomrule
\end{tabular}
\end{center}

\vspace{2ex}

Memorylessness is the thread through the waiting-time column. The Geometric and the Exponential are the unique discrete and continuous distributions with that property, which is why both admit one-line recursive expectation arguments.
\end{keybox}
\newpage


\subsubsection{Select Worked Problems}

\begin{problem}
A basketball player makes 70\% of free throws. In 10 attempts, what's $P(\text{exactly 8 makes})$?
\end{problem}

\begin{solution}
We are given a target out of 10 indepdent trials, each with the same probability, so we can use the Binomial distribution:
\[
P(X = 8) = \binom{10}{8} \cdot 0.7^8 \cdot 0.3^2 \approx 0.233
\]
\end{solution}

\vspace{-2ex}

\begin{problem}
A website gets 2 signups per hour on average. What's $P(\text{0 signups in the next hour})$?
\end{problem}

\begin{solution}
The problem gives a rate over a timeframe, so this is uses the Poisson or Exponential, this time we can use either. Working in hours, the rate is $\lambda = 2$ per hour and the window is $t = 1$.

\textbf{Poisson.} Count the signups directly and ask for zero of them:
\[
P(X = 0) = e^{-\lambda}\frac{\lambda^0}{0!} = e^{-2} \approx 0.135.
\]

\textbf{Exponential.} Zero signups in the hour is the same event as the wait for the first signup exceeding one hour, which is the tail probability $P(X > t) = 1 - F(t)$:
\[
P(X > 1) = 1 - \left(1 - e^{-\lambda t}\right) = e^{-\lambda t} = e^{-2} \approx 0.135.
\]

The two agree, which is the identity $P(N_t = 0) = P(X > t)$ between a Poisson count of zero and an Exponential wait past $t$.
\end{solution}


\begin{problem}
You keep drawing cards (with replacement) from a deck until you get an ace. Expected number of draws?
\end{problem}

\begin{solution}
We want to know the expected trials until the first success with replacement, which is the mean of the Geometric distrbution:

\[
\E[X] = \mu = \frac{1}{p} = \frac{1}{4/52} = 13
\]
\end{solution}

\vspace{-2ex}

\begin{problem}
A factory produces items with a 5\% defect rate. In a batch of 20, what's the expected number of defective items and $P(\text{at most 1 defective})$?
\end{problem}

\begin{solution}
We want the successes (defects) in 20 fixed trials, which calls for the Binomial distribution:
\[
\E[X] = \mu = np = 20 \cdot 0.05 = 1
\]

\[
P(X \leq 1) = P(0) + P(1) = \left ( \binom{20}{0} \cdot 0.05^0 \cdot 0.95^{20} \right ) + \left ( \binom{20}{1} \cdot 0.05^1 \cdot 0.95^{19} \right ) = 0.358 + 0.377 = 0.736
\]
\end{solution}

\newpage

\begin{problem}
On average, a bus arrives every 10 minutes. What's $P(\text{waiting more than 15 minutes})$?
\end{problem}

\begin{solution}
We want the tail probability of the Exponential distribution, making our unit of time 1 minute, then $\lambda = 0.1$. 
\[
P(X > 15) = 1 - (1 - e^{-0.1 \cdot 15}) = e^{-1.5} \approx 0.223
\]
\end{solution}

\vspace{-2ex}

\begin{problem}
Test scores are Normal with mean 70 and standard deviation 10. What fraction of students score above 90?
\end{problem}

\begin{solution}
This is simply asking what proportion of a normal distribution exists in the upper right tail two standard deviations from the mean. We know that 95.4\% of the data lies within 2$\sigma$, then the portion only in the right tail is $\frac{1-0.954}{2} = 0.023 = 2.3\%$
\end{solution}

\vspace{-1ex}

\begin{problem}
You roll a die until you've gotten two 6's total. Expected number of rolls?
\end{problem}

\begin{solution}
We can use linearity of expectation along with the Geometric distribution (or more directly, the Negative Binomial). Then
\[
\E[X] = \frac{r}{p} = \frac{2}{1/6} = 12
\]
\end{solution}

\vspace{-2ex}

\begin{problem}
Errors on a page occur at a rate of 0.5 per page (Poisson). A book has 200 pages. What's the expected total errors? What's the approximate probability the book has fewer than 90 errors?
\end{problem}

\begin{solution}
Scaling the rate to the full book, the total is Poisson with
\[
\E[X] = 200 \cdot 0.5 = 100 = \lambda, \qquad \sigma = \sqrt{\lambda} = 10.
\]
The book total is a sum of 200 independent per-page Poisson counts, so by the CLT it is approximately Normal, and with $\lambda = 100$ well above the $\lambda \approx 10$ threshold the approximation is safe: $X \approx \mathcal{N}(100, 100)$.

Fewer than 90 errors is one standard deviation below the mean, so by the 68\% rule
\[
P(X < 90) \approx \frac{1 - 0.68}{2} = 0.16.
\]
\end{solution}

\vspace{-2ex}

\begin{problem}
You flip a coin with P(H) = p. What's the variance of the number of heads in n flips? For what value of p is the variance maximized?
\end{problem}

\begin{solution}
This is a binary outcome with a fixed number of flips, which calls for the Binomial distribution. The variance is therefore $\sigma^2 = np(1-p)$. This is maximized at $p = \frac{1}{2}$.
\end{solution}


\newpage

% ============================================================
\section{Brainteasers}
% ============================================================

% ============================================================
\subsection{Strategies}
% ============================================================

Most of these problems are solved by recognizing a recurring structure rather than by extended calculation. A given problem usually reduces to one of a small set of standard techniques, and the central skill is identifying which one applies. The approach proceeds in three phases: frame the problem precisely, map it to a known reduction, and verify the result against limiting cases. 

% ============================================================
\subsubsection{Framing}
% ============================================================

Much of the work happens at the setup stage. An error introduced while framing a problem propagates through every later step and cannot be repaired by correct algebra, so the setup warrants explicit attention before any technique is chosen.

\begin{keybox}[Framing checklist]
\begin{itemize}[itemsep=1ex, label=--]
  \item State every implicit assumption: fairness of coins and dice, independence of trials, sampling with or without replacement, and the exact meaning of ``random'', meaning which distribution over which set
  \item Fix the sample space and the unit of observation. Counting the wrong objects as equally likely is the most common source of probability errors
  \item Classify the target quantity as a probability, an expectation, a count, a ratio, an optimal strategy, or a worst-case guarantee. The class selects the toolset
\end{itemize}
\end{keybox}

The choice of observational unit deserves particular care. A single experiment supports several sample spaces, and only some of them make the outcomes equally likely. A direct count is valid only at the level where that symmetry holds, so the first task is to identify that level and count there.

% ============================================================
\subsubsection{Core Probabilistic Reductions}
% ============================================================

Most probability puzzles reduce to one of the techniques below:

\vspace{1ex}

\begin{keybox}[Cue to reduction]
\begin{center}
\small
\setlength{\tabcolsep}{5pt}
\begin{tabular}{@{}ll@{}}
\toprule
\textbf{Cue in the problem} & \textbf{Reduction} \\
\midrule
``at least one'' & complement, $1 - P(\text{none})$ \\
``expected number of \dots'' & linearity of expectation with indicators \\
exchangeable items, ``$A$ before $B$'' & symmetry \\
``until the first \dots'', repeated identical trials & conditioning on the first step \\
waiting for a pattern, random walk, absorbing states & Markov states \\
updating on evidence, false positives & Bayes, odds form \\
fixed trial count, rate, or waiting structure & recognize the named distribution \\
sum or average of many independent pieces & normal approximation, CLT \\
\bottomrule
\end{tabular}
\end{center}
\end{keybox}

Three techniques recur often enough to warrant restating:

\textbf{Linearity with indicators.} The expected size of a count equals the sum of the probabilities of its parts, whether or not those parts are independent. The method is to write the count as a sum of indicator variables, one per event, and take expectations term by term. Its power comes from sidestepping dependence entirely, since linearity holds without any independence assumption, which reduces an intractable joint distribution to a list of individual probabilities.

\textbf{Symmetry.} When several items are exchangeable, every statement about their relative order is uniform over the possible orderings. This converts a probability into a counting ratio with no further work. Two consequences appear repeatedly: a specific item $A$ precedes a specific item $B$ with probability $\tfrac{1}{2}$, and among $k$ competing items each is equally likely to be first. When the items instead compete at rates $\lambda_i$, the uniform $\tfrac{1}{k}$ generalizes to the rate share $\lambda_j / \sum_i \lambda_i$.

\textbf{Conditioning and self-similarity.} A process is self-similar when its state after one step is a scaled copy of the original problem. Conditioning on that first step then expresses the unknown in terms of itself, producing a solvable equation. This is recursive expectation for waiting times and first-step analysis for Markov chains, and it is the standard route for any ``expected time until'' quantity.

% ============================================================
\subsubsection{Combinatorial and Logical Tactics}
% ============================================================

These tactics address counting problems and questions of the form ``is a configuration reachable'' or ``must something exist''.

\textbf{Small cases and induction.} Compute the answer for $n = 1, 2, 3$, conjecture the pattern, then confirm it by induction or a direct argument. Beyond producing the formula, small cases expose the recursive structure that conditioning later exploits.

\textbf{Invariants.} An invariant is a quantity left unchanged by every legal move. If two states differ in the value of some invariant, no sequence of moves connects them, which resolves reachability and ``is it possible'' questions without any search. The usual construction assigns each configuration a value in a small set, often through a parity or coloring argument, and shows that every move preserves it.

\textbf{Pigeonhole and extremal arguments.} If more items than categories are distributed among the categories, some category holds at least two. The extremal variant examines the largest or smallest element of a configuration to force a contradiction or an existence claim. Both establish that something must occur without exhibiting it, which is what makes them suited to existence questions.

\textbf{Working backwards.} When the goal state is simple and the starting state is complex, reason from the end. The value of the final position is known outright, and each earlier position is defined in terms of positions closer to the goal, which is the structure underlying optimal-stopping and game problems.

\textbf{Generating functions.} Encoding a sequence as the coefficients of a power series turns operations on the sequence into algebra on the series. Convolution becomes multiplication, so counting problems with additive constraints and distributions of sums of independent variables both reduce to expanding a product.

% ============================================================
\subsubsection{Adversarial and Sequential Decisions}
% ============================================================

These problems involve a decision maker choosing actions over time, sometimes against an opponent, to optimize an outcome.

\textbf{Backward induction.} In a finite game of alternating moves with perfect information, solve from the last move outward. At each position a player selects the move that optimizes their outcome assuming optimal play thereafter. Applied recursively, this fixes the value of every position together with the equilibrium strategy.

\textbf{Dominant strategies and minimax.} A dominant strategy is optimal regardless of the opponent's choice, and when one exists the analysis reduces to it. When none exists, a zero-sum game is solved by minimax, in which each player minimizes their maximum loss. With mixed strategies permitted, the two players' values coincide at a single value of the game, which is von Neumann's minimax theorem.

\textbf{Optimal stopping.} When observations arrive in sequence and each accept-or-reject decision is final, the problem is one of optimal stopping, and its solution is typically a threshold rule. Observe without committing for a calibrated initial portion of the sequence to establish a baseline, then accept the first later observation that exceeds it. In the canonical case of selecting the best of $n$ candidates presented in random order, the initial portion is a fraction $1/e$ and the resulting success probability also approaches
\[
\frac{1}{e} \approx 0.368
\]

\textbf{Fair games and martingales.} A fair bet has expected value zero, and a sequence of fair bets forms a martingale, whose expected future value equals its present value regardless of history. No stopping rule applied with finite capital converts a fair game into a favorable one in expectation, which is the formal reason betting systems fail. Gambler's Ruin is the standard model, and its ruin and duration formulas quantify the outcome.

\begin{remark}[Kelly sizing]
For a favorable repeated bet, the wealth fraction that maximizes long-run growth is the edge divided by the odds, which for an even-money bet at win probability $p$ is $2p - 1$. Staking more than this fraction raises variance without raising growth, and staking well beyond it drives the long-run growth rate negative despite the positive edge.
\end{remark}

% ============================================================
\subsubsection{Counting Information}
% ============================================================

A separate class of problems asks for the minimum number of measurements that guarantee a result in the worst case. The governing tool is outcome counting. A measurement with $b$ distinguishable results partitions the remaining possibilities into $b$ groups, so $k$ measurements distinguish at most $b^k$ cases. Resolving $M$ possibilities therefore requires
\[
k \geq \log_b M
\]

\begin{keybox}[Worst-case guarantees by outcome counting]
\begin{center}
% \footnotesize
\begin{tabular}{@{}lcl@{}}
\toprule
\textbf{Measurement} & \textbf{Outcomes per step} & \textbf{$k$ steps resolve} \\
\midrule
yes/no question & $2$ & $2^k$ \\
balance weighing & $3$ & $3^k$ \\
\bottomrule
\end{tabular}
\end{center}

\vspace{2ex}

A balance produces three outcomes, left heavy, right heavy, or level, so $k$ weighings resolve $3^k$ cases, and locating one heavy item among $N$ of known direction needs $\lceil \log_3 N \rceil$. If the odd item may be either heavier or lighter, each candidate splits into two hypotheses and the requirement becomes $3^k \geq 2N$.
\end{keybox}

The same bound gives $\lceil \log_2 M \rceil$ yes/no questions, achieved by binary search. When the possibilities are not equally likely, the bound sharpens from a raw count to the entropy: the expected number of yes/no questions needed is at least $H(X)$ bits.

% ============================================================
\subsubsection{Estimation}
% ============================================================

Order-of-magnitude questions are handled by decomposition. Break the unknown into a product of factors, each estimable to within a small multiple, anchor on a known reference quantity, estimate each factor independently, and combine. Because the factor errors are multiplicative and roughly independent, they tend to partially cancel rather than accumulate, so a chain of rough estimates usually lands within a factor of two or three of the true value.

% ============================================================
\subsubsection{Verification}
% ============================================================

A closed-form answer should survive a set of inexpensive checks before it is trusted. Each check tests the answer against a case where the correct value is already known or constrained.

\vspace{1.5ex}

\begin{keybox}[Sanity checks]
\begin{itemize}[itemsep=1ex, label=--]
  \item \textbf{Boundary.} Send a parameter to an extreme, such as $p \to 0$, $p \to 1$, $n = 1$, or $n \to \infty$, and confirm the formula reduces to the obvious value
  \item \textbf{Range.} A probability lies in $[0,1]$, an expected number of trials is at least $1$, and a variance is nonnegative
  \item \textbf{Monotonicity.} The answer should move in the correct direction as each parameter increases
  \item \textbf{Symmetry.} If the setup is symmetric in two variables, the answer must be symmetric in them as well
  \item \textbf{Small case.} Evaluate by hand at $n = 2$ or $3$ and match the result against the formula
\end{itemize}
\end{keybox}
\newpage

% ============================================================
\subsection{Select Worked Problems}
% ============================================================

\begin{problem}[Easy]
Five pirates need to divide 100 Gold Coins. Pirates have a hierarchy, from Level 1 to level 5. The highest level pirate is the leader. The leader proposes a plan to distribute the gold and all the pirates vote on it (including the leader). If at least 50\% of the pirates agree on the plan, the gold is split according to the proposal. If not, level-5 pirate is kicked from the ship, and the level-4 pirate now proposes a new plan. This process continues until a proposal is accepted. All pirates are extremely smart and extremely greedy. How does level-5 Pirate divide the treasure?
\end{problem}

\begin{solution}
We approach this problem game-theoretically and dynamically:
\begin{itemize}[label=$\circ$]
    \item 2 pirates $\rightarrow$ Pirate 2 proposes 100 coins for themself and zero for Pirate 1
    \item 3 pirates $\rightarrow$ Pirate 3 bribes Pirate 1 with one coin and keeps 99 for themself, which they accept because they get zero coins if Pirate 3 is booted off.
    \item 4 pirates $\rightarrow$ Pirate 4 bribes Pirate 2 with one coin and keeps 99 for themself, in the same manner as the 3-pirate case.
    \item 5 pirates $\rightarrow$ Pirate 5 bribes pirates 1 and 3 with one coin each, keeping 98 for themself.
\end{itemize}

\begin{center}
Thus, the final division is [98, 0, 1, 0 ,1].
\end{center}

\end{solution}


\begin{problem}[Easy]
An explorer travels 1 mile South, 1 mile East, then 1 mile North, and finds themselves back at the starting point. The explorer spots a bear, what is the color of the bear?
\end{problem}

\begin{solution}
The bear must be white. There are a few relevant scenarios that could have resulted in the explorer starting back where they began. In the first case, they started at the North pole, then travelling South equates to travelling directly away from the Pole. They can then travel \textit{any} distance East or West, the following mile back North will return them to their destination. There are no bears in the South Pole so we omit this, though the same logic applies (with the Eastward movement equivalent to spinning in place or not moving at all).

There are other solutions near the South Pole. If the traveler's initial one-mile journey South ends on a latitude circle with a circumference of $\frac{1}{n}\  (n \in \mathbb{Z}^+)$ miles, walking one mile East will complete $n$ full rotations around the pole and leave them exactly where they started walking East, meaning the following mile North would place them in the same starting position.
\end{solution}

\begin{problem}[Easy]
There is a light bulb inside a room and three switches outside. All switches are currently in off state and only one switch controls the light bulb. You may turn any number of switches on or off any number of times you want. How many times do you need to go into the room to figure out which switch controls the light bulb?
\end{problem}

\begin{solution}
We only need to enter the room once. The first switch is flipped and left on for enough time that the bulb become hot. It is then flipped off and the second switch is flipped, we then immediately enter the room. This leads the three cases:
\begin{enumerate}
    \item Bulb is hot $\rightarrow$ the light must be off and the first switch controls the light
    \item Bulb is on $\rightarrow$ the second switch controls the light
    \item Bulb is off and cold $\rightarrow$ the third switch controls the light
\end{enumerate}

\medskip

If there were four switches, we can still infer the switch in one entry by turning on both switches 1 and 2 for a long time, then switching off switch 2 and turning on switch 3 before entering the room. This results in the following cases:
\begin{enumerate}
    \item Bulb is on and hot $\rightarrow$ Switch 1 controls the light
    \item Bulb is on and cool $\rightarrow$ Switch 3 controls the light
    \item Bulb is off and hot $\rightarrow$ Switch 2 controls the light
    \item Bulb is off and cool $\rightarrow$ Switch 4 controls the light
\end{enumerate}
\end{solution}

The extension to four bulbs worked because we had two binary variables (illumination and bulb temperature) which creates 4 dinstinct states, of which we were only using 3 in the 3-bulb case.

\medskip

\begin{problem}[Easy]
How do you place 50 good candies and 50 rotten candies in two boxes such that if you choose a box at random and take out a candy at random, it better be good!

We need to maximize the probability of getting a good candy when selecting a random box and a random candy from it.
\end{problem}

\begin{solution}
On one extreme, we can place all of the good candies in one box with all of the bad candies in another. This gives us $\mathbb{E}[G] = 0.5$. Thus, we should find solutions that better than 50\%. We can beat this rather simply, by placing one good candy in one box and the other 49 in the other box with all of the bad candies. This results in:
\[
\mathbb{E}[G] = 0.5 * \frac{1}{1} + 0.5 * \frac{49}{99} \approx 0.747
\]
\end{solution}

\begin{problem}[Easy]
In a world where everyone wants a girl child, each family continues having babies till they have a girl. What do you think will the boy-to-girl ratio be eventually? 
\end{problem}

\begin{solution}
The stopping rule shapes each family but does not affect the aggregate ratio, since it cannot bias any individual birth. Each birth is an independent fair coin, and the choice of when to stop having children cannot change the overall fraction of each sex.

Counting round by round makes this precise. Begin with $N$ couples. In the first round all $N$ give birth, producing $\tfrac{N}{2}$ girls and $\tfrac{N}{2}$ boys, and the couples with a girl stop. The $\tfrac{N}{2}$ couples with a boy have another child, producing $\tfrac{N}{4}$ girls and $\tfrac{N}{4}$ boys, and the process repeats. Each round is an independent batch of fair births and contributes equal numbers of girls and boys, so the totals remain equal and the ratio is
\[
1 : 1
\]

The appearance of a girl surplus comes from the family structure. Every family terminates on exactly one girl, which makes girls salient as the final child. The boys are the births preceding that girl, and their count follows a geometric distribution with $p = \tfrac{1}{2}$ and mean $\tfrac{1-p}{p} = 1$. The average family therefore has one boy and one girl. The rule fixes the position of the girl within each family, not the proportion of each sex overall.
\end{solution}

\begin{problem}
Hundred tigers and one sheep are put on a magic island that only has grass. Tigers can live on grass, but they want to eat sheep. If a Tiger bites the Sheep then it will become a sheep itself. If 2 tigers attack a sheep, only the first tiger to bite converts into a sheep. Tigers don’t mind being a sheep, but they have a risk of getting eaten by another tiger. All tigers are intelligent and want to survive. Will the sheep survive?
\end{problem}

\begin{solution}
This is a parity argument by backward induction, sharing the alternating logic of the pirate loot-splitting problem. The cases build from the bottom up, and the governing fact is that a tiger eats the sheep only when doing so leaves it safe, meaning the resulting one-sheep island has an outcome in which that new sheep is not itself eaten.

Consider the smallest cases. With one tiger, no other tiger threatens it, so it eats the sheep and survives. With two tigers, a tiger that eats the sheep becomes a sheep facing one remaining tiger, which is the one-tiger case in which the sheep is eaten, so neither tiger risks it and the sheep survives. With three tigers, a tiger that eats reduces the island to the two-tiger case, in which a sheep survives, so eating is safe and the sheep is eaten. With four, eating reduces to the three-tiger case, in which the sheep is eaten, so no tiger risks it and the sheep survives.

Each case is thus determined by the one below it: the sheep is eaten precisely when the case with one fewer tiger leaves the sheep safe. This alternation follows parity, so the sheep survives for an even number of tigers and is eaten for an odd number. With 100 tigers, the sheep survives.
\end{solution}

\begin{problem}
A duck is sitting at the center of a circular lake. A fox is waiting at the shore, not able to swim, wishing to eat the duck. The Fox can move around the whole lake at a speed four times the speed at which the duck can swim. The duck can fly, but only once it reaches the shore of the lake, it can't fly from the water directly. Can the duck always reach the shore without being eaten by the fox?
\end{problem}

\begin{solution}
Let $R$ be the radius of the lake, $v$ the duck's speed, and $4v$ the fox's. A direct dash across fails: the duck would swim $R$ while the fox runs at most half the circumference, $\pi R$. Since $\pi < 4$, the fox covers its distance in less time and arrives first.

Instead the duck swims in a small circle concentric with the lake, of radius slightly less than $\tfrac{R}{4}$. The choice of $\tfrac{R}{4}$ comes from comparing angular speeds. A duck at radius $r$ sweeps angle at rate $\tfrac{v}{r}$, while the fox running the shore sweeps at rate $\tfrac{4v}{R}$, so the duck out-rotates the fox exactly when
\[
\frac{v}{r} > \frac{4v}{R} \iff r < \frac{R}{4}
\]
The fox tracks the duck from the nearest point of the shore, but since the duck gains angle steadily, the separation between them grows until the duck sits diametrically opposite the fox.

At that moment the duck turns and swims straight for the shore. It must cover $R - \tfrac{R}{4} = \tfrac{3R}{4}$, while the fox must run half the circumference, $\pi R$. Their times are
\[
t_{\text{duck}} = \frac{3R}{4v}, \qquad t_{\text{fox}} = \frac{\pi R}{4v}
\]
so the comparison reduces to $3 < \pi$, which holds. The duck reaches the shore first and escapes.
\end{solution}

\newpage

\begin{problem}
An airplane flies in a straight line from airport $A$ to airport $B$, and then back in a straight line from $B$ to $A$. It travels with a constant engine speed and there is no wind. Will the total travel time for the same round trip be greater, less or the same if, throughout both the flights a constant wind blows from $A$ to $B$?
\end{problem}


\begin{solution}
The round trip takes longer with wind. Take a plane speed of $100$ mph, a wind of $10$ mph, and a distance of $100$ miles each way. The wind adds to the ground speed on one leg and subtracts on the other, giving
\[
\frac{100}{110} + \frac{100}{90} \approx 2.02 \text{ hours}
\]
against $2$ hours with no wind.

In general, for distance $d$, plane speed $v$, and wind $w < v$, the round trip takes
\[
\frac{d}{v+w} + \frac{d}{v-w} = \frac{2dv}{v^2 - w^2}
\]
which exceeds the windless time $\frac{2d}{v}$ for any $w \neq 0$, since $v^2 - w^2 < v^2$. The asymmetry is that the plane spends less time on the fast leg and more time on the slow one, so the time lost to the headwind always outweighs the time gained from the tailwind.
\end{solution}

\begin{problem}
You are given two cords that both burn in exactly one hour, not necessarily at a uniform rate. How should you light the cords to determine a time interval of exactly 15 minutes?
\end{problem}

\begin{solution}
We light three of the four ends at the same time. When the cord with both ends lit burns out, thirty minutes will have passed. At this point we snuff the fire on the other cord. The segment remaining can be lit from both sides and will burn for an interval of 15 minutes.

The same interval can be obtained from a single cord by maintaining four simultaneous flames, assuming that cutting and lighting take negligible time. A cord holds sixty minutes of burn time, and with $k$ flames alight that reserve is consumed at rate $k$, so four flames exhaust it in $\frac{60}{4} = 15 \text{ minutes}$

To begin, we cut the cord in two and light all four ends. Whenever a piece burns out and the count drops to two flames, we cut one of the remaining pieces and light the two new ends, restoring the count to four. The number of cuts required is not bounded in advance, since it depends on how unevenly the cord burns, but the flame count is held at four throughout and the cord is spent after exactly fifteen minutes.
\end{solution}

\begin{problem}
There are 10 black socks and 10 red socks in a cupboard. Your task is to draw a few socks at random, without looking, such that you get at least one matching pair (single color). What's the least number of socks that you need to draw?
\end{problem}

\begin{solution}
You only need to draw three. If the first two were opposite colours (not a match), then the third sock \textit{must} match one of the first two. This is the pigeonhole principle with 3 pigeons (socks drawn) and two holes (colours).
\end{solution}

\newpage

\begin{problem}
Prove that there are two opposite points on the Earth's surface that have the same temperature. Assume that the temperature varies continuously over the surface.
\end{problem}

\begin{solution}
Consider any great circle on the surface, such as the equator, and let $t(\theta)$ be the temperature at angle $\theta$. Two points are opposite when their angles differ by $\pi$, so the goal is a value of $\theta$ satisfying $t(\theta) = t(\theta + \pi)$. Define $f(\theta) = t(\theta) - t(\theta + \pi)$, which is zero when the two points have equal temperatures. As a difference of continuous functions, $f$ is continuous.

Evaluating at the ends of a half circle, and using $t(2\pi) = t(0)$ because the circle closes on itself,
\[
f(0) = t(0) - t(\pi), \qquad f(\pi) = t(\pi) - t(0) = -f(0)
\]
so $f$ takes values of opposite sign at $0$ and $\pi$, or is zero, in which case the temperatures are equal and we are done.Otherwise, by the Intermediate Value Theorem there must be a point on the closed interval $[0, \pi]$ such that $f(\theta) = 0$, and that root is an antipodal pair at equal temperature.

This is the one-dimensional case of the Borsuk-Ulam theorem, which states that any continuous map $S^n \rightarrow \mathbb{R}^n$ admits an antipodal pair with equal values. Thematically it resembles the hairy ball theorem, in that a continuous structure on a sphere is forced into a degeneracy, though Borsuk-Ulam constrains values at antipodal pairs while the hairy ball theorem forces a tangent vector field to vanish somewhere.
\end{solution}

\begin{problem}
The probability of having at least one accident on a road in one hour is $\frac{3}{4}$. What is the probability that at least one accident occurs in half an hour?
\end{problem}

\begin{solution}
Assuming accidents in disjoint intervals are independent, no accident in an hour means no accident in each of its two half-hour intervals. Let $q$ be the probability of no accident in half an hour. Then $q^2 = 1 - \frac{3}{4} = \frac{1}{4} \rightarrow q = \frac{1}{2}$. Then the probability of at least one accident in half an hour is $1-q = \frac{1}{2}$.
\end{solution}

\begin{problem}
Suppose you're on a game show, and you're given the choice of three doors: Behind one door is a car; behind the others, goats. You pick a door, say No. 1, and the host, who knows what's behind the doors, opens another door, say No. 3, which has a goat. Now, do you want to pick door No. 2? What is the probability to win the car if you switch?
\end{problem}

\begin{solution}
Switching wins with probability $\frac{2}{3}$.

Suppose door 1 is chosen. The car is equally likely to be behind each door, giving three cases. If the car is behind door 1, the host opens one of the goats and switching loses. If the car is behind door 2, the host is forced to open door 3, and switching wins. If the car is behind door 3, he is forced to open door 2, and switching wins.

Switching therefore wins in two of the three equally likely cases. Equivalently, switching wins exactly when the initial pick was wrong, which occurs with probability $\frac{2}{3}$.
\end{solution}

\newpage

\begin{problem}
One hundred prisoners are lined up, facing one direction and assigned a random hat, either red or blue. Each prisoner can see the hats in front of them but not behind. Starting with the prisoner at the back of the line and moving forward, they must each, in turn, say only one word which must be "red" or "blue". If the word matches their hat color they are released, if not, they are kept imprisoned. They can hear each others' answers, no matter how far they are on the line, but they do not hear the verdict (whether the answer was correct). A friendly guard warns them one night before, giving them enough time to come up with a strategy. How many prisoners can be freed using the best strategy? Assume that there is an unknown number of red \& blue hats.
\end{problem}

\begin{solution}
Ninety-nine prisoners are guaranteed release, and the first to speak is correct with probability $\frac{1}{2}$. This is optimal, since that prisoner sees only the hats ahead and receives no information bearing on their own, so no strategy improves on a coin flip for them.

The prisoners agree in advance to encode parity. The prisoner at the back, who sees all 99 hats ahead, says ``red'' if the number of red hats ahead is even and ``blue'' if it is odd.

Each later prisoner deduces their own colour from three known quantities: the announced parity, the colours called behind them, and the hats visible ahead. Their hat is red exactly when the combined parity of the reds heard and seen differs from the announced parity, since their own hat is the only one unaccounted for and so must absorb the discrepancy. Having deduced it, they answer correctly, and their answer joins the record used by those ahead.
\end{solution}

\begin{problem}
A line of 100 airline passengers is waiting to board a plane. They each hold a ticket to one of the 100 seats on that flight. For convenience, let's say that the $n_{\text{th}}$ passenger in line has a ticket for seat number `$n$'. Being drunk, the first person in line picks a random seat (equally likely for each seat). All of the other passengers are sober, and will go to their assigned seats unless it is already occupied; If it is occupied, they will then find a free seat to sit in, at random. What is the probability that the last (100th) person to board the plane will sit in their own seat (\#100)?
\end{problem}

\begin{solution}
Label seats so passenger $n$ owns seat $n$, with passenger 1 drunk.

First we pin down where the 100th passenger can end up. Take any passenger $k$ with $2 \le k \le 99$. When $k$ boards, seat $k$ is either free, in which case $k$ takes it, or already occupied. Either way seat $k$ is occupied once $k$ has boarded. So when passenger 100 boards, seats $2,\dots,99$ are full and the one free seat is 1 or 100.

Now we show these are equally likely. Seats 1 and 100 are the only seats no passenger in $\{2,\dots,99\}$ owns. Both begin free, and every passenger who chooses at random chooses uniformly among the free seats. Swap the labels 1 and 100. Passenger 1 chooses at random no matter what he owns, passenger 100 takes the leftover no matter what she owns, and the owners of seats $2,\dots,99$ are unaffected. So the swap sends the boarding process to an equally likely one with the two seats exchanged, and the two outcomes have equal probability.

So, passenger 100 gets seat 100 with probability $\tfrac12$.
\end{solution}

\begin{problem}
A stick is broken into 3 parts, by choosing 2 points randomly along its length. With what probability can it form a triangle?
\end{problem}

\begin{solution}
Call the length $\ell$ and set $\ell = 1$. To form a triangle every piece must
be shorter than $\tfrac12$, since a piece of length $\tfrac12$ or more equals or
exceeds the sum of the other two.

Fix the first cut and measure it from its nearer end, calling that distance $X$,
so $X$ is uniform on $[0, \tfrac12]$. Let $Y$ be the second cut, uniform on
$[0,1]$.

If $Y < X$ the pieces are $Y$, $X - Y$, and $1 - X$. Here $1 - X \ge \tfrac12$, so
this always fails.

If $Y > X$ the pieces are $X$, $Y - X$, and $1 - Y$. We already have $X < \tfrac12$.
The remaining piece $1 - Y$ stays below $\tfrac12$ when $Y > \tfrac12$, and the
middle piece $Y - X$ stays below $\tfrac12$ when $Y < X + \tfrac12$. So we succeed
exactly when
\[
  \tfrac12 < Y < X + \tfrac12
\]
a window of width $X$. Given $X$, the success probability is therefore $X$.

Since the first cut is uniform on $[0, \tfrac12]$, we can average with $\frac{1}{b-a} \int_a^b g(x)dx$,
\[
  P(\text{triangle}) = \frac{1}{\frac{1}{2} - 0}\int_0^{1/2} x\, dx
  = 2\left[\frac{x^2}{2}\right]_0^{1/2} = \left(\tfrac12\right)^2 - 0 = \tfrac14
\]
So the probability of forming a valid triangle is $\tfrac14$.
\end{solution}

\begin{problem}
A rabbit sits at the bottom of a staircase with $n$ stairs. The rabbit can hop up only one or two stairs at a time. What kind of sequence is depicted by the different ways possible for the rabbit to ascend to the top of the stairs of length $n=1,2,3\dots$?
\end{problem}

\begin{solution}
We'll call the number of paths to a certain stair $n$ as $f(n)$. We can compute $f(1)$ to $f(5)$ as:

\begin{center}
\begin{tabular}{c c l}
\hline
$n$ & $f(n)$ & paths \\
\hline
1 & 1 & \texttt{1} \\
2 & 2 & \texttt{11, 2} \\
3 & 3 & \texttt{111, 12, 21} \\
4 & 5 & \texttt{1111, 112, 121, 211, 22} \\
5 & 8 & \texttt{11111, 1112, 1121, 1211, 2111, 122, 212, 221} \\
\hline
\end{tabular}
\end{center}

\vspace{1em}

Examining each level from $n=3$ onwards, we can see that each path can be constructed as the paths to $n-2$ plus a double hop, or the paths to $n-1$ plus a single hop. Thus, the number of paths to $n$ is equal to the number of paths to $n-2$ plus the number of paths to $n-1$, which is the Fibonacci recurrence. So the pattern that arises is the Fibonacci
sequence.

\end{solution}


\begin{problem}
$A, B$, and $C$ live together and share everything equally. One day $A$ brings home 5 logs of wood, $B$ brings 3 logs and $C$ brings none. Then they use the wood to cook together and share the food. Since $C$ did not bring any wood, he gives \$8 instead. How much to $A$ and how much to $B$?
\end{problem}

\begin{solution}
The three of them consume $8$ logs total, so each eats $\frac{8}{3}$ logs.
Measuring what each contributed beyond his own share, $A$ has a surplus of
$5 - \frac{8}{3} = \frac{7}{3}$ logs and $B$ has a surplus of
$3 - \frac{8}{3} = \frac{1}{3}$ logs.

$C$'s \$8 pays for the $\frac{8}{3}$ logs he ate without providing any. Those
logs were supplied entirely out of the two surpluses, so the money splits in
proportion to surplus:
\[
  \frac{7}{3} : \frac{1}{3} = 7 : 1
\]
Thus $A$ receives \$7 and $B$ receives \$1.
\end{solution}

\newpage

\begin{problem}
\vspace{-1ex}
Suppose you have a hotel which has one floor with infinite number of rooms in a row and all of them are occupied. A new customer wants to check in, how will you accommodate her? What if infinite number of people want to check in, how will you accommodate them? 

Lastly, suppose infinite number of buses arrive at the hotel, each having infinite number of people, how will you accommodate them?

\end{problem}

\begin{solution}
This is a question of counting infinities, or more specifically, of finding a 1-1 mapping that includes the new incoming people. Since all of these are based on integer amounts of people, each is countable (even an infinite number of buses of infinite people).

In the first case, we can simply ask each person in room $k$ to move to room $k+1$. Similarly for the second case, we can ask each $k_{\text{th}}$ person to move to room $2 \cdot k$. This would leave open every $2 \cdot k - 1$ room open for the infinite number of people arriving on the bus.

For the final case, identify each person by a pair. An original occupant is
labeled by their room number $k \ge 1$. A bus passenger is labeled $(i, j)$,
meaning seat $j \ge 1$ of bus $i \ge 1$. We assign rooms in successive shells,
indexed by $n = 1, 2, 3, \dots$, where shell $n$ contains:
\begin{itemize}
  \item the $n$th original occupant,
  \item seats $1$ through $n$ of bus $n$, that is the pairs $(n, 1), \dots, (n, n)$,
  \item seat $n$ of every earlier bus, that is the pairs $(1, n), \dots, (n-1, n)$.
\end{itemize}

The bus pairs in shell $n$ are exactly those with $\max(i, j) = n$, so every pair
$(i, j)$ lands in shell $\max(i, j)$ and lands in no other. Every occupant $k$
lands in shell $k$. Each shell holds $2n$ people, a finite count, so processing
the shells in order $1, 2, 3, \dots$ reaches every person after finitely many
rooms are filled. No bus is skipped, since a fixed pair $(i, j)$ is served by the
time we finish shell $\max(i, j)$.
\end{solution}

\begin{problem}
\vspace{-1ex}
A group has 70 members. For any two members $X$ and $Y$, there is a language
that $X$ speaks but $Y$ does not, and there is a language that $Y$ speaks but
$X$ does not. At least how many different languages are spoken by this group?
\end{problem}

\begin{solution}
Start by considering a smaller group, say $n = 6$ with members labelled 1
to 6. We show this is supportable by 4 languages, denoted $A$, $B$, $C$, and
$D$. Assign languages as sets:
\begin{align*}
1 &= \{A, B\} & 2 &= \{A, C\} & 3 &= \{A, D\} \\
4 &= \{B, C\} & 5 &= \{B, D\} & 6 &= \{C, D\}
\end{align*}

\vspace{-0.5em}

The condition is that for any two members $X$ and $Y$, each speaks a language the
other does not. In set terms, neither set may contain the other. Here all six
sets have size 2, and no size-2 set can contain another (collections of distinct but differently sized sets could contain subsets), so any two distinct sets
satisfy the condition automatically.

This suggests giving every member a set of the same size $r$, which is both
convenient and optimal. Convenient, because the number of available sets is then just $\binom{n}{r}$. Optimal, because nothing is gained by mixing sizes: enlarging one member's set
only shrinks the pool of sets that can coexist with it. By Sperner's theorem, the
largest antichain (a collection of sets where none contains another) built from
$n$ elements is obtained by taking all subsets of a single size
$\lfloor n/2 \rfloor$, giving $\binom{n}{\lfloor n/2 \rfloor}$ of them. 

Finally, the smallest $n$ for which $\binom{n}{\lfloor n/2 \rfloor} \ge 70$ is $n=8$ languages. 
\end{solution}


\newpage

\begin{problem}
Two players, $A$ and $B$, form a team playing against an opponent. To enter,
the team pays \$1 at the start. If the team wins it receives \$3, and nothing
otherwise.

$A$ and $B$ are placed in two separate rooms with no communication between them.
Each flips a fair coin. $A$ must guess the outcome of $B$'s flip, and $B$ must
guess the outcome of $A$'s flip. The team wins only if both guesses are correct,
and wins nothing if either guess is wrong. The two may agree on a strategy in
advance, before being separated.

Should the team pay \$1 to play?
\end{problem}

\begin{solution}
The two coins land in one of four equally likely states:
\[
  \begin{array}{cccc}
    \{H,H\} & \{H,T\} & \{T,H\} & \{T,T\}
  \end{array}
\]
With no strategy, each player guesses the other coin at random. Each guess is
correct with probability $\tfrac12$, and the two guesses are independent, so both
are correct with probability $\tfrac12 \cdot \tfrac12 = \tfrac14$. The expected
return is $\mathbb{E}[\text{Play}] = \tfrac14 \cdot 3 + \tfrac34 \cdot 0 - 1 = -\tfrac14$, which is negative, so random guessing is not worth playing. 

Instead, the strategy is for each player to guess that the other coin \textit{equals their
own}. Under this rule, $A$'s guess is correct exactly when $B$'s coin equals
$A$'s, and $B$'s guess is correct exactly when $A$'s coin equals $B$'s. These are
the same condition: both coins show the same face. So the two guesses are no
longer independent, they are both correct together or both wrong together.

The team therefore wins precisely when the coins match. Of the four states, two
match and two do not:
\[
  \underbrace{\{H,H\}, \; \{T,T\}}_{\text{match, win}}
  \qquad
  \underbrace{\{H,T\}, \; \{T,H\}}_{\text{no match, lose}}
\]
so the win probability is $\tfrac24 = \tfrac12$. The expected return becomes
\[
  \mathbb{E}[\text{Play}] = \tfrac12 \cdot 3 + \tfrac12 \cdot 0 - 1 = \tfrac12 \quad \longrightarrow \quad \text{which is positive, so they should play.}
\]
\end{solution}

\begin{problem}
Snow particles fall to the ground one after another. Each snowflake is of type
``Stellar Dendrite'' with probability $p$ if the preceding particle was also a
Stellar Dendrite, and with probability $q$ if the preceding particle was some
other type. If a snowflake is picked at random from the ground, what is the
probability that it is a Stellar Dendrite?
\end{problem}

\begin{solution}
Let $x$ be the probability that a snowflake picked up is an SD. We set this up as
a Markov chain and use first-step analysis, conditioning on the previous flake.

Assume the system has reached a steady state, so the probability of an SD does not
change from one flake to the next. Then the previous flake is an SD with the same
probability $x$. The current flake is an SD in two ways: the previous was an SD
(probability $x$) and it continues (probability $p$), or the previous was not
(probability $1 - x$) and it starts (probability $q$). Since the current flake is
itself a picked-up flake, its SD probability is also $x$:
\[
  x = x \cdot p + (1 - x) \cdot q
\]
Solving,
\[
  x = xp + q - qx \implies x - xp + qx = q \implies x = \frac{q}{1 - p + q}
\]
\end{solution}

\newpage

\begin{problem}
A drawer contains red socks and black socks. When two socks are drawn
at random, the probability that both are red is $\frac{1}{2}$. (a) How small can the number of socks in the drawer be? (b) How small if the number of black
socks is even? 
\end{problem}

\begin{solution}
Draws are without replacement, so for red socks $r$ and black socks $b$ the probability that both are red factors as
\[
p = \frac{r}{r+b} \times \frac{r-1}{r+b-1} = \frac{1}{2}
\]
We can test this with $r=3$ and $b=1$:
\[
\frac{3}{4} \times \frac{2}{3} = \frac{1}{2}
\]
As desired. For (b), the second factor is smaller than the first when $b > 0$, so $\frac{r}{r+b} > \frac{r-1}{r+b-1}$. Multiplying this inequality by each factor and using $p = \frac{1}{2}$ gives
\[
\left( \frac{r}{r+b} \right)^2 > \frac{1}{2} > \left( \frac{r-1}{r+b-1} \right)^2
\]
We solve these for bounds on $r$ in terms of $b$, then check the admissible candidate $r$ for each even $b$ until the product equals $\frac{1}{2}$.
\end{solution}

\begin{problem}
Four people, $A$, $B$, $C$ and $D$ need to get across a river. The only way to cross the river is
by an old bridge, which holds at most 2 people at a time. Being dark, they can't cross the
bridge without a torch, of which they only have one. So each pair can only walk at the
speed of the slower person. They need to get all of them across to the other side as
quickly as possible. $A$ is the slowest and takes 10 minutes to cross; $B$ takes 5 minutes; $C$
takes 2 minutes; and $D$ takes 1 minute.
What is the minimum time to get all of them across to the other side?
\end{problem}

\begin{solution}
We could first try using $D$ to `ferry' everyone across, coming back each time for the next person, but this takes 19 minutes and never uses people on \textit{both} sides of the bridge.

Instead, send $C$ and $D$ over (2 min), $D$ back (3 min), $A$ and $B$ over (13 min), $C$ back (15 min), then $C$ and $D$ over again for 17 minutes total. The trick is to pair the two slowest so we pay $A$'s 10 minutes on the same trip as $B$'s, leaving a fast runner across to carry the torch back.

This beats ferrying whenever $2t_2 < t_1 + t_3$, where the $t_i$ are the sorted times: two crossings by the second fastest must cost less than one each by the fastest and second slowest. Here $2C < D + B$, or $4 < 6$, saving 2 minutes.
\end{solution}

\newpage

\begin{problem}
Boss $A$'s birthday is one of the following 10 dates:
\begin{itemize}
    \item Mar 4, Mar 5, Mar 8
    \item Jun 4, Jun 7
    \item Sep 1, Sep 5
    \item Dec 1, Dec 2, Dec 8
\end{itemize}
$A$ tells you only the month and tells colleague $C$ only the day.
\begin{enumerate}
    \item You state: ``I do not know $A$'s birthday, and $C$ does not know it either.''
    \item $C$ replies: ``I did not know $A$'s birthday, but now I do.''
    \item You state: ``Now I know it, too.''
\end{enumerate}
An assistant overhears the exchange and correctly determines the birthday. What is $A$'s birthday?
\end{problem}

\begin{solution}
Since we know the month, our personal list of candidate birthdays is restricted to that single month. For our colleague $C$, the same principle applies with respect to the day.

From (1), we announce that we do not know the birthday and that $C$ cannot possibly know it either. This means our month cannot contain any unique day numbers (7 in June, 2 in December), which eliminates both June and December entirely. The candidate list is reduced to the March and September dates:
\begin{itemize}
    \item March: 4, 5, 8
    \item September: 1, 5
\end{itemize}

From (2), $C$ announces that they now know the birthday. This rules out day 5, since 5 appears in both March and September and would leave $C$ uncertain. The remaining candidate dates are March 4, March 8, and September 1.

From (3), we announce that we now know the birthday as well. If our month were March, we would still have two possibilities (March 4 and March 8). Because we can determine the exact date, our month must be September, making the boss's birthday September 1.
\end{solution}

\begin{problem}
A casino offers a card game using a normal deck of 52 cards. The rule is that you turn
over two cards each time. For each pair, if both are black, they go to the dealer's pile; if
both are red, they go to your pile; if one black and one red, they are discarded. The
process is repeated until you two go through all 52 cards. If you have more cards in your
pile, you win \$100; otherwise (including ties) you get nothing. The casino allows you to
negotiate the price you want to pay for the game. How much would you be willing to
pay to play this game?
\end{problem}

\begin{solution}
The objective is to have more cards in our pile than the house, so we want some way to ``burn'' black cards. However, the only way that black cards (and symmetrically, red cards) can be used up without awarding points to the house is if they are discarded. Note, however, that this requires a red card burn too. Thus, there is \textit{no} way to build an advantage. In fact, every game will always end in a tie (though the amount of discarded and held cards may differ game to game). The game is unwinnable, so no amount of money makes it worth playing. 
\end{solution}
\newpage

% ============================================================
\section{Code}
% ============================================================

% ============================================================
% SQL Patterns (PostgreSQL) -- master wrapper for sections/30-code/31-sql/.
% Mirrors the sections/30-code.tex <-> sections/30-code/ pattern one
% level down. \input paths are relative to the MAIN document being
% compiled, so paths below are given in full from the project root.
% ============================================================
\subsection{SQL Patterns (PostgreSQL)}
% ============================================================

Reference for recurring query shapes. Each entry states the problem phrasing
that signals it and the mechanism that makes it work.

\begin{keybox}[Reading a problem]
Most prompts map to a small set of shapes, and the wording is the tell:
\begin{itemize}[label=$\circ$, itemsep=1ex]
  \item ``per group'', ``for each'' $\to$ \texttt{GROUP BY}, or a
  \texttt{PARTITION BY} window when the whole row must survive.
  \item ``top-$N$ per group'', ``most recent per'' $\to$ rank-then-filter,
  or \texttt{DISTINCT ON} when $N = 1$.
  \item ``change from the previous'', ``period-over-period'', ``return''
  $\to$ \texttt{LAG} / \texttt{LEAD}.
  \item ``running'', ``cumulative'', ``rolling'', ``moving average'' $\to$
  an ordinary aggregate with \texttt{OVER}.
  \item ``consecutive'', ``streak'', ``gap'' $\to$ gaps and islands.
  \item ``one column per category'', ``pivot'' $\to$ conditional aggregation.
\end{itemize}

\vspace{1ex}
The first decision on any prompt is whether the answer collapses rows
(\texttt{GROUP BY}) or keeps every row and appends a computed column (a
window function).
\end{keybox}

\begin{keybox}[Logical processing order]
A query is evaluated in this order, regardless of written order:
\[
  \texttt{FROM/JOIN} \to \texttt{WHERE} \to \texttt{GROUP BY} \to
  \texttt{HAVING} \to \texttt{SELECT} \to \texttt{DISTINCT} \to
  \texttt{ORDER BY} \to \texttt{LIMIT}
\]
Two consequences drive most of the patterns below. \texttt{WHERE} runs
before aggregation, so it cannot see aggregates or \texttt{SELECT} aliases.
Window functions are computed at the \texttt{SELECT} step, so they also
cannot appear in \texttt{WHERE} and must be wrapped in a subquery or CTE
to filter on.
\end{keybox}

\subsubsection{Filtering and Joins}

The foundation. Every later pattern assumes you can already restrict rows and
combine tables.

\paragraph{WHERE vs HAVING.} \texttt{WHERE} filters individual rows before
grouping. \texttt{HAVING} filters whole groups after aggregation and is the
only place an aggregate predicate is legal. A filter on a plain
(non-aggregated) column belongs in \texttt{WHERE} so rows are dropped before
the grouping work happens. Reserve \texttt{HAVING} for conditions like
\texttt{COUNT(*) > 5}.

\begin{lstlisting}[style=sql, caption={WHERE prunes rows, HAVING prunes groups}]
SELECT department, COUNT(*) AS headcount
FROM employees
WHERE active = true          -- row filter, before grouping
GROUP BY department
HAVING COUNT(*) > 5;         -- group filter, on an aggregate
\end{lstlisting}

\vspace{-1ex}

\newpage

\paragraph{Joins.} The recognition cue is which side you want to keep when a
match is missing.
\begin{itemize}[label=$\circ$]
  \item \textbf{INNER:} rows with a match on both sides.
  \item \textbf{LEFT / RIGHT:} keep all rows from one side, \texttt{NULL}-fill the other.
  \item \textbf{FULL OUTER:} keep unmatched rows from both sides.
  \item \textbf{CROSS:} Cartesian product, every left row against every right row.
  \item \textbf{SELF:} a table joined to itself under an alias (employee to manager, row to prior row).
  \item \textbf{Semi-join} (\texttt{EXISTS}, \texttt{IN}): keep left rows with at least one match. No row duplication, no right-side columns.
  \item \textbf{Anti-join} (\texttt{NOT EXISTS}, or \texttt{LEFT JOIN ... WHERE right.key IS NULL}): keep left rows with no match.
\end{itemize}

Semi- and anti-joins filter the left table on whether a match exists, without
pulling right-side columns or duplicating rows. Prefer \texttt{NOT EXISTS} over
\texttt{NOT IN} for the anti-join, since a \texttt{NULL} on the right poisons
\texttt{NOT IN} to zero rows.

\begin{lstlisting}[style=sql, caption={Semi-join keeps matches, anti-join keeps non-matches}]
-- Semi-join: customers with at least one order
SELECT c.customer_id, c.name
FROM customers c
WHERE EXISTS (SELECT 1 FROM orders o
              WHERE o.customer_id = c.customer_id);

-- Anti-join: customers with no orders
SELECT c.customer_id, c.name
FROM customers c
WHERE NOT EXISTS (SELECT 1 FROM orders o
                  WHERE o.customer_id = c.customer_id);

-- Anti-join, LEFT JOIN form: unmatched rows are the NULL-filled ones
SELECT c.customer_id, c.name
FROM customers c
LEFT JOIN orders o ON o.customer_id = c.customer_id
WHERE o.customer_id IS NULL;
\end{lstlisting}
\vspace{-4ex}

\subsubsection{Aggregation}

\paragraph{Conditional aggregation (pivot).} The signal is a request for one
column per category on a single row. Fold a categorical column into counts by
wrapping a \texttt{CASE} in an aggregate. Rows the \texttt{CASE} does not match
contribute 0 to \texttt{SUM}. In Postgres the \texttt{FILTER} clause does the
same thing and reads more cleanly.

\begin{lstlisting}[style=sql, caption={Pivot risk categories to columns}]
SELECT
    inspection_type,
    COUNT(*) FILTER (WHERE risk_category = 'Low Risk')      AS low_risk,
    COUNT(*) FILTER (WHERE risk_category = 'High Risk')     AS high_risk,
    COUNT(*)                                                AS total
FROM sf_restaurant_health_violations
GROUP BY inspection_type
ORDER BY total DESC;
\end{lstlisting}

\vspace{-1ex}

The \texttt{SUM(CASE WHEN ... THEN 1 ELSE 0 END)} form is equivalent and
portable across engines. \texttt{COUNT} ignores \texttt{NULL}, so
\texttt{COUNT(CASE WHEN cond THEN 1 END)} (no \texttt{ELSE}) also counts only
the matches.

\subsubsection{NULL Handling}

\texttt{NULL} means unknown, so it propagates through arithmetic and poisons
comparisons. The two functions below are the standard guards.

\paragraph{COALESCE and NULLIF.} \texttt{COALESCE(a, b, c, ...)} returns the
first non-\texttt{NULL} argument, evaluated left to right, or \texttt{NULL} if
all are \texttt{NULL}. The common use is a fallback that keeps a \texttt{NULL}
from poisoning arithmetic: \texttt{salary + commission} is \texttt{NULL}
whenever \texttt{commission} is, so write \texttt{salary + COALESCE(commission, 0)}.
It also supplies a default down a priority list, as in
\texttt{COALESCE(nickname, first\_name, 'Unknown')}. It is ANSI standard,
unlike the vendor-specific \texttt{IFNULL} (MySQL) and \texttt{ISNULL}
(SQL Server), and it accepts any number of arguments.

\texttt{NULLIF(a, b)} is the mirror image: it returns \texttt{NULL} when
\texttt{a = b}, otherwise \texttt{a}. Its standard use guards division,
\texttt{x / NULLIF(y, 0)} yields \texttt{NULL} instead of raising a
divide-by-zero when \texttt{y} is 0.

\begin{lstlisting}[style=sql, caption={NULLIF guards division, COALESCE supplies defaults}]
SELECT
    order_id,
    revenue / NULLIF(quantity, 0)  AS unit_price,  -- NULL, not an error, when quantity = 0
    COALESCE(discount, 0)          AS discount      -- NULL discount treated as 0
FROM orders;
\end{lstlisting}

\vspace{-1ex}

A subtlety with aggregates: \texttt{SUM}, \texttt{AVG}, and
\texttt{COUNT(col)} skip \texttt{NULL}s. So \texttt{AVG(commission)} averages
only the non-\texttt{NULL} rows, while \texttt{AVG(COALESCE(commission, 0))}
counts the \texttt{NULL}s as zeros and enlarges the denominator. 

\paragraph{Gotchas.}
\begin{itemize}
  \item \textbf{\texttt{NOT IN} with NULLs:} if the subquery returns any
  \texttt{NULL}, \texttt{x NOT IN (...)} is never true and the query returns
  zero rows. \texttt{x NOT IN (1, NULL)} expands to
  \texttt{NOT (x=1 OR x=NULL)}, and the \texttt{NULL} comparison poisons the
  result to unknown. Use \texttt{NOT EXISTS}, which is null-safe.
  \item \textbf{Integer division truncates:} \texttt{5 / 2} yields \texttt{2}
  because the fractional part is discarded. Cast one operand:
  \texttt{5::numeric / 2} or \texttt{5.0 / 2}. Bites hardest in ratio and
  percentage columns.
  \item \textbf{\texttt{COUNT(col)} vs \texttt{COUNT(*)}:} the former skips
  \texttt{NULL}s in \texttt{col}, the latter counts every row. A common source
  of off-by-$N$ errors.
\end{itemize}

\subsubsection{Window Functions}

The signal for a window function is a per-row result that still depends on
other rows, such as a rank, a neighbor value, or a running total. Where
\texttt{GROUP BY} would collapse the group to one row, a window function keeps
every row and appends the computed column.

Every window function shares the same shape:
\texttt{f() OVER (PARTITION BY ... ORDER BY ...)}. \texttt{PARTITION BY} splits
the rows into independent groups, \texttt{ORDER BY} sets the order within each
group, and the function is evaluated per row. All of them are computed at the
\texttt{SELECT} step, so to filter on their output you wrap them in a subquery
or CTE.

\paragraph{Rank-then-filter with ROW\_NUMBER().} Use this for top-$N$ per
group when you need to keep the full row, since \texttt{GROUP BY} would collapse
the other columns away. Number rows inside each partition, then filter on that
number in an outer query.

Choose the ranker by tie policy: \texttt{ROW\_NUMBER} assigns distinct numbers
and breaks ties arbitrarily, \texttt{RANK} repeats a rank and leaves gaps,
\texttt{DENSE\_RANK} repeats with no gaps.

\begin{lstlisting}[style=sql, caption={Highest-paid employee per department}]
SELECT department, name, salary
FROM (
    SELECT department, name, salary,
           ROW_NUMBER() OVER (PARTITION BY department
                              ORDER BY salary DESC) AS rn
    FROM employees
) ranked
WHERE rn = 1;
\end{lstlisting}

\vspace{-4ex}

\paragraph{Offset access with LAG() and LEAD().} The signal is a comparison
against the previous or next row: a delta, a period-over-period change, or a
return. These read a value from another row relative to the current one,
without a self-join. \texttt{LAG(col, k, default)} returns the value \texttt{k}
rows earlier in the ordered partition, \texttt{LEAD} looks \texttt{k} rows
ahead. Both default to \texttt{k = 1} and return \texttt{NULL} past the
partition edge unless a fill \texttt{default} is supplied. The \texttt{ORDER BY}
inside \texttt{OVER} defines what ``previous'' means and is mandatory.
\texttt{PARTITION BY} resets the sequence at each group boundary, so a lag never
bleeds from one partition into the next.

\begin{lstlisting}[style=sql, caption={Daily returns per ticker with LAG}]
SELECT ticker, trade_date, close,
       close / LAG(close) OVER (PARTITION BY ticker
                                ORDER BY trade_date) - 1 AS daily_return
FROM prices;
\end{lstlisting}

\vspace{-1ex}

The first row of each partition has no predecessor, so its \texttt{LAG} and the
derived return are \texttt{NULL}, which is correct since there is no prior price
to compare against. \texttt{LEAD} is the same tool aimed forward, used below for
gap detection.

\paragraph{Running and rolling aggregates.} The signal is ``running'',
``cumulative'', or ``moving average''. An ordinary aggregate gains a window when
you attach \texttt{OVER (... ORDER BY ...)}. With an \texttt{ORDER BY} and no
explicit frame, the frame defaults to everything from the partition start up to
the current row, giving a running total. Add an explicit \texttt{ROWS BETWEEN}
frame to get a fixed moving window instead.

\begin{lstlisting}[style=sql, caption={20-day moving average and running volume}]
SELECT ticker, trade_date, close,
       AVG(close) OVER (PARTITION BY ticker ORDER BY trade_date
                        ROWS BETWEEN 19 PRECEDING AND CURRENT ROW) AS ma_20,
       SUM(volume) OVER (PARTITION BY ticker ORDER BY trade_date)  AS cum_volume
FROM prices;
\end{lstlisting}

\vspace{-1ex}

\texttt{ma\_20} is the current row plus the 19 before it, and \texttt{cum\_volume}
uses the default frame for a running total. These are the SQL forms of a pandas
\texttt{.rolling(20).mean()} and \texttt{.expanding().sum()}.

One subtlety: the default frame is \texttt{RANGE UNBOUNDED PRECEDING}, which
includes every row tied on the \texttt{ORDER BY} key, not just rows physically
up to the current one. When ties on the order key exist and you want a
row-precise window, state \texttt{ROWS BETWEEN} explicitly.

\paragraph{Gaps and islands.} The signal is any mention of consecutive values,
streaks, or runs. An \emph{island} is a maximal run of consecutive values. The
trick: subtract a dense row number from the value. Across a consecutive run
both sides increase by one per step, so the difference stays constant, and it
jumps exactly at a gap. Group by the difference to collapse each run.

\begin{lstlisting}[style=sql, caption={Islands of consecutive ids}]
WITH numbered AS (
    -- DISTINCT first so duplicate ids do not desync the row number
    SELECT DISTINCT log_id,
           ROW_NUMBER() OVER (ORDER BY log_id) AS rn
    FROM logs
)
SELECT MIN(log_id) AS island_start,
       MAX(log_id) AS island_end,
       COUNT(*)    AS length
FROM numbered
GROUP BY log_id - rn
ORDER BY island_start;
\end{lstlisting}

The \texttt{DISTINCT} is load-bearing. If \texttt{log\_id} repeats, \texttt{rn}
still advances on the duplicate while the value does not, so
\texttt{log\_id - rn} shifts and splits an island that should stay whole. For
consecutive \emph{dates} rather than integers, offset by an interval:
\texttt{log\_date - rn * INTERVAL `1 day'} is constant within a run.

\begin{lstlisting}[style=sql, caption={Gaps via LEAD}]
SELECT log_id + 1  AS missing_from,
       next_id - 1 AS missing_to
FROM (
    SELECT log_id,
           LEAD(log_id) OVER (ORDER BY log_id) AS next_id
    FROM logs
) t
WHERE next_id - log_id > 1;   -- values missing between two present ids
\end{lstlisting}

\subsubsection{Deduplication}

\paragraph{DISTINCT.} \texttt{DISTINCT} compares the entire select list as a
tuple. \texttt{SELECT DISTINCT a, b} returns every distinct \emph{pair}, so a
value in \texttt{a} can still repeat as long as the paired \texttt{b} differs.
It is a whole-row operator with no way to say ``dedupe on this column and keep
the rest of the row.'' It runs after \texttt{SELECT} and before
\texttt{ORDER BY}, which is why every \texttt{ORDER BY} expression must also
appear in the select list, and it treats two \texttt{NULL}s as duplicates.
\texttt{COUNT(DISTINCT col)} is a separate mechanism that dedupes the input to
one aggregate and has no effect on the rows returned.

\newpage

\paragraph{DISTINCT ON.} A PostgreSQL extension that keeps the first row of
each group sharing the same value of the parenthesized expressions, returning
all selected columns of that surviving row. ``First'' is defined by
\texttt{ORDER BY}, whose leading terms must match the \texttt{DISTINCT ON}
expressions. Omit the \texttt{ORDER BY} and the surviving row is arbitrary.

\begin{lstlisting}[style=sql, caption={Most recent close per ticker}]
SELECT DISTINCT ON (ticker)
       ticker, trade_date, close
FROM prices
ORDER BY ticker, trade_date DESC;   -- leading term must be ticker
\end{lstlisting}

\vspace{-1ex}

This is the compact form of the rank-then-filter pattern from the window
section. The \texttt{ROW\_NUMBER} equivalent needs a subquery and an outer
\texttt{WHERE rn = 1}. Which to reach for:

\begin{itemize}[label=$\circ$]
  \item \textbf{Top-1 per group in Postgres:} \texttt{DISTINCT ON} is shorter
  and generally at least as fast, since it stops at the first row of each
  group.
  \item \textbf{Top-$N$ for $N > 1$, or portability:} \texttt{ROW\_NUMBER}
  extends to \texttt{rn <= 3} and runs on any engine. \texttt{DISTINCT ON} is
  Postgres-only.
  \item \textbf{Ties:} \texttt{DISTINCT ON} always emits exactly one row per
  group, settled by the remaining \texttt{ORDER BY} terms or arbitrarily. To
  keep every tied row, use \texttt{RANK() = 1}.
\end{itemize}

Because the \texttt{ORDER BY} is dictated by the deduplication key, presenting
the result in a different order means wrapping the query in an outer
\texttt{SELECT ... ORDER BY}.

\subsubsection{String Manipulation}

Positions are 1-based, and every function returns a new value rather than
mutating in place. The recurring uses are normalizing a dirty join key, parsing
a delimited field, and folding a group of strings into one.

\paragraph{Concatenation.} The operator \texttt{a || b} propagates
\texttt{NULL}: if either side is \texttt{NULL} the whole expression is
\texttt{NULL}. \texttt{CONCAT(a, b, ...)} treats \texttt{NULL} as an empty
string. \texttt{CONCAT\_WS(sep, ...)} takes the separator first, joins the rest
with it, and skips \texttt{NULL} arguments without leaving a doubled separator.
Choosing \texttt{||} for user-facing text is a common way to silently blank out
a whole column.

\paragraph{Case and whitespace.} \texttt{UPPER} and \texttt{LOWER} are the
standard folds, \texttt{INITCAP} capitalizes the first letter of each word and
lowercases the rest. \texttt{TRIM} strips characters from the ends only, spaces
by default, with \texttt{LEADING}, \texttt{TRAILING}, or \texttt{BOTH} to pick
a side. Interior whitespace survives, so collapsing double spaces needs
\texttt{REGEXP\_REPLACE(s, '\textbackslash s+', ' ', 'g')}.
\texttt{LPAD(s, n, fill)} and \texttt{RPAD} pad to a fixed width and truncate
when the input is already longer than \texttt{n}.

\vspace{2ex}

\begin{lstlisting}[style=sql, caption={Normalizing a join key}]
SELECT
    LOWER(TRIM(email))                      AS email_key,
    SPLIT_PART(LOWER(TRIM(email)), '@', 2)  AS domain,
    LPAD(account_id::text, 8, '0')          AS padded_id,
    CONCAT_WS(', ', last_name, first_name)  AS sort_name
FROM users;
\end{lstlisting}

\newpage

\paragraph{Extraction and search.} \texttt{SUBSTRING(s FROM a FOR b)} (or
\texttt{SUBSTR(s, a, b)}) takes \texttt{b} characters starting at position
\texttt{a}. \texttt{LEFT(s, n)} and \texttt{RIGHT(s, n)} take from the ends,
and a negative \texttt{n} drops that many characters from the opposite end.
\texttt{STRPOS(s, sub)} returns the 1-based index of the first occurrence and
\texttt{0} when the substring is absent, so a found/not-found test compares
against \texttt{0}. \texttt{SPLIT\_PART(s, delim, n)} pulls the \texttt{n}th
field of a delimited string and returns an empty string when that field does
not exist, worth wrapping in \texttt{NULLIF(..., '')}.

\begin{lstlisting}[style=sql, caption={Parsing a delimited name field}]
SELECT full_name,
       SPLIT_PART(full_name, ' ', 1)              AS first_name,
       NULLIF(SPLIT_PART(full_name, ' ', 2), '')  AS last_name,
       LENGTH(full_name)
         - LENGTH(REPLACE(full_name, ' ', ''))    AS space_count
FROM people;
\end{lstlisting}

\vspace{-1ex}

The last expression is the occurrence-counting idiom: remove every copy of the
target and measure how much the string shrank. For a needle longer than one
character, divide the difference by \texttt{LENGTH(needle)}.

\paragraph{Replacement.} \texttt{REPLACE(s, from, to)} substitutes every
literal occurrence. \texttt{TRANSLATE(s, from, to)} maps character to character
by position and deletes any character in \texttt{from} that has no counterpart
in \texttt{to}, a fast way to strip a known character set.
\texttt{REGEXP\_REPLACE(s, pattern, repl, flags)} handles the general case and
replaces only the \emph{first} match unless the \texttt{`g'} flag is passed, a
frequent source of half-cleaned output.

\paragraph{Pattern matching.} \texttt{LIKE} uses \texttt{\%} for any run of
characters and \texttt{\_} for exactly one, and it is case sensitive.
\texttt{ILIKE} is the case-insensitive form. For real expressiveness the POSIX
operators apply: \texttt{\textasciitilde} matches, \texttt{\textasciitilde*}
matches case-insensitively, and \texttt{!\textasciitilde} negates. Anchors
\texttt{\^{}} and \texttt{\$} matter, since an unanchored pattern matches
anywhere in the string.

\begin{lstlisting}[style=sql, caption={Regex cleaning and validation}]
SELECT phone,
       REGEXP_REPLACE(phone, '[^0-9]', '', 'g')  AS digits,
       phone ~ '^\+?[0-9]{10,15}$'               AS looks_valid,
       email ILIKE '%@gmail.com'                 AS is_gmail
FROM contacts;
\end{lstlisting}

\vspace{-1ex}

\paragraph{Splitting and aggregating.} \texttt{STRING\_AGG(expr, delim ORDER BY
...)} collapses a group into one delimited string and is the string analogue of
\texttt{SUM}. The \texttt{ORDER BY} lives inside the call, so the output order
is deterministic. The inverse direction expands one row into many:
\texttt{UNNEST(STRING\_TO\_ARRAY(s, delim))} or
\texttt{REGEXP\_SPLIT\_TO\_TABLE(s, pattern)} in the \texttt{FROM} clause.

\begin{lstlisting}[style=sql, caption={Collapse a group, then expand it back}]
SELECT department,
       STRING_AGG(name, ', ' ORDER BY name) AS members
FROM employees
GROUP BY department;
-- one row per tag stored in a comma-separated column
SELECT p.id, TRIM(tag) AS tag
FROM products p,
     UNNEST(STRING_TO_ARRAY(p.tags, ',')) AS tag;
\end{lstlisting}


\newpage

\subsubsection{Command Reference}

Signatures are PostgreSQL. Optional arguments are shown in brackets.

\begin{table}[H]
\centering
\footnotesize
\renewcommand{\arraystretch}{1.4}
\rowcolors{2}{gray!10}{white}
\begin{tabularx}{\textwidth}{>{\raggedright\arraybackslash}p{0.20\textwidth} >{\raggedright\arraybackslash}p{0.225\textwidth} >{\raggedright\arraybackslash}X}
\toprule
\textit{Command} & \textit{Arguments} & \textit{What it does} \\
\midrule
\textbf{WHERE} & boolean predicate & Filters individual rows before grouping. Cannot see aggregates, window functions, or select aliases. \\
\textbf{GROUP BY} & column or expression list & Collapses rows to one output row per distinct key combination. \\
\textbf{HAVING} & boolean predicate & Filters whole groups after aggregation. The only legal home for an aggregate predicate. \\
\textbf{DISTINCT} & applies to the select list & Removes duplicate rows compared across the entire select list. Treats two \texttt{NULL}s as equal. \\
\textbf{DISTINCT ON} & \texttt{(expr, ...)} & Postgres. Keeps one row per distinct value of the expressions, chosen by the leading \texttt{ORDER BY} terms, which must match. \\
\textbf{ORDER BY} & \texttt{expr [ASC|DESC] [NULLS FIRST|LAST]} & Sorts the output. Postgres places \texttt{NULL}s last under \texttt{ASC} and first under \texttt{DESC}. \\
\textbf{LIMIT / OFFSET} & \texttt{n} & Caps the number of rows returned, or skips the first \texttt{n}. \\
\textbf{WITH} & \texttt{name AS (query)} & Common table expression. Names a subquery for reuse. \texttt{WITH RECURSIVE} walks hierarchies. \\
\textbf{CASE} & \texttt{WHEN cond THEN val [ELSE val] END} & Row-level branching. Returns the first match, or \texttt{NULL} when nothing matches and \texttt{ELSE} is absent. \\
\textbf{FILTER} & \texttt{(WHERE cond)} after an aggregate & Restricts which rows the aggregate consumes. Postgres shorthand for \texttt{SUM(CASE ...)}. \\
\textbf{EXISTS / NOT EXISTS} & subquery & Semi-join and anti-join. Null-safe, unlike \texttt{IN} and \texttt{NOT IN}. \\
\bottomrule
\end{tabularx}
\caption*{Query clauses and structure}
\end{table}

\vspace{-2ex}

\begin{table}[H]
\centering
\footnotesize
\renewcommand{\arraystretch}{1.4}
\rowcolors{2}{gray!10}{white}
\begin{tabularx}{\textwidth}{>{\raggedright\arraybackslash}p{0.20\textwidth} >{\raggedright\arraybackslash}p{0.225\textwidth} >{\raggedright\arraybackslash}X}
\toprule
\textit{Command} & \textit{Arguments} & \textit{What it does} \\
\midrule
\textbf{COUNT(*)} & none & Counts every row, including all-\texttt{NULL} ones. \\
\textbf{COUNT(col)} & one column & Counts rows where \texttt{col} is non-\texttt{NULL}. \\
\textbf{COUNT(\allowbreak DISTINCT col)} & one column & Counts distinct non-\texttt{NULL} values. \\
\textbf{SUM / AVG /\allowbreak MIN / MAX} & one expression & Skip \texttt{NULL}s. \texttt{AVG} divides by the non-\texttt{NULL} count. \\
\textbf{STRING\_AGG} & \texttt{(expr, delim [ORDER BY ...])} & Concatenates a group into one delimited string. \\
\textbf{ARRAY\_AGG} & \texttt{(expr [ORDER BY ...])} & Collects group values into an array. \\
\textbf{COALESCE} & any number of values & Returns the first non-\texttt{NULL} argument, left to right. \\
\textbf{NULLIF} & two values & Returns \texttt{NULL} when the arguments are equal, otherwise the first. Guards division by zero. \\
\textbf{GREATEST / LEAST} & list of values & Largest or smallest value across columns within one row. Ignores \texttt{NULL}s. \\
\textbf{IS [NOT] DISTINCT FROM} & two values & Null-safe comparison that treats two \texttt{NULL}s as equal. \\
\textbf{IS [NOT] NULL} & one value & The only valid null test. \texttt{= NULL} is always unknown. \\
\bottomrule
\end{tabularx}
\caption*{Aggregates and \texttt{NULL} handling}
\end{table}

\vspace{-2ex}

\begin{table}[H]
\centering
\footnotesize
\renewcommand{\arraystretch}{1.4}
\rowcolors{2}{gray!10}{white}
\begin{tabularx}{\textwidth}{>{\raggedright\arraybackslash}p{0.20\textwidth} >{\raggedright\arraybackslash}p{0.225\textwidth} >{\raggedright\arraybackslash}X}
\toprule
\textit{Command} & \textit{Arguments} & \textit{What it does} \\
\midrule
\textbf{OVER} & \texttt{(PARTITION BY ... ORDER BY ... [frame])} & Turns a function into a window function evaluated per row without collapsing the result set. \\
\textbf{ROW\_NUMBER()} & none & Sequential number within the partition. Ties broken arbitrarily. \\
\textbf{RANK()} & none & Competition rank. Ties share a rank and leave gaps afterward. \\
\textbf{DENSE\_RANK()} & none & Ties share a rank with no gaps. \\
\textbf{NTILE(n)} & bucket count & Splits the partition into \texttt{n} roughly equal buckets and labels each row. \\
\textbf{LAG(col[, k, default])} & \texttt{k} defaults to 1 & Value \texttt{k} rows earlier in the ordered partition. \texttt{NULL} past the partition edge unless a default is supplied. \\
\textbf{LEAD(col[, k, default])} & same as \texttt{LAG} & Value \texttt{k} rows ahead. \\
\textbf{FIRST\_VALUE /\allowbreak LAST\_VALUE} & \texttt{col} & First or last value in the current frame. \texttt{LAST\_VALUE} needs an explicit frame to see the full partition. \\
\textbf{NTH\_VALUE(col, n)} & column, position & The \texttt{n}th value in the frame. \\
\textbf{SUM / AVG / COUNT ... OVER} & aggregate plus window spec & Running or rolling aggregate computed per row. \\
\textbf{ROWS BETWEEN} & \texttt{a PRECEDING} to \texttt{b FOLLOWING}, \texttt{CURRENT ROW}, \texttt{UNBOUNDED} & Physical frame counted in rows. Use it for fixed moving windows. \\
\textbf{RANGE BETWEEN} & same bound keywords & Value-based frame that includes every row tied on the \texttt{ORDER BY} key. This is the default frame. \\
\bottomrule
\end{tabularx}
\caption*{Window functions and frames}
\end{table}

\newpage

\begin{table}[H]
\centering
\footnotesize
\renewcommand{\arraystretch}{1.4}
\rowcolors{2}{gray!10}{white}
\begin{tabularx}{\textwidth}{>{\raggedright\arraybackslash}p{0.20\textwidth} >{\raggedright\arraybackslash}p{0.225\textwidth} >{\raggedright\arraybackslash}X}
\toprule
\textit{Command} & \textit{Arguments} & \textit{What it does} \\
\midrule
\textbf{\texttt{||}} & two strings & Concatenation. Yields \texttt{NULL} if either side is \texttt{NULL}. \\
\textbf{CONCAT} & any number of values & Concatenation treating \texttt{NULL} as an empty string. \\
\textbf{CONCAT\_WS} & separator, then values & Joins arguments with a separator and skips \texttt{NULL}s. \\
\textbf{LENGTH / CHAR\_LENGTH} & one string & Character count. \texttt{OCTET\_LENGTH} gives bytes. \\
\textbf{UPPER / LOWER /\allowbreak INITCAP} & one string & Case folding. \texttt{INITCAP} capitalizes each word and lowercases the rest. \\
\textbf{TRIM} & \texttt{[LEADING|TRAILING|BOTH] [chars] FROM s} & Strips characters from the ends only, spaces by default. \texttt{LTRIM}, \texttt{RTRIM}, \texttt{BTRIM} are shorthands. \\
\textbf{LPAD / RPAD} & \texttt{(s, n [, fill])} & Pads to width \texttt{n}, truncating when \texttt{s} is already longer. \\
\textbf{LEFT / RIGHT} & \texttt{(s, n)} & First or last \texttt{n} characters. A negative \texttt{n} drops that many from the opposite end. \\
\textbf{SUBSTRING} & \texttt{(s FROM a FOR b)} or \texttt{SUBSTR(s, a, b)} & Slice of length \texttt{b} starting at 1-based position \texttt{a}. \\
\textbf{STRPOS / POSITION} & \texttt{(s, sub)} / \texttt{(sub IN s)} & 1-based index of the first occurrence, \texttt{0} when absent. \\
\textbf{REPLACE} & \texttt{(s, from, to)} & Replaces every literal occurrence. \\
\textbf{TRANSLATE} & \texttt{(s, from\_set, to\_set)} & Maps characters by position and deletes those with no counterpart. \\
\textbf{SPLIT\_PART} & \texttt{(s, delim, n)} & The \texttt{n}th delimited field. Returns an empty string when the field is missing. \\
\textbf{STRING\_TO\_\allowbreak ARRAY / UNNEST} & \texttt{(s, delim)} / \texttt{(array)} & Splits into an array, then expands an array into rows. \\
\textbf{REVERSE / REPEAT} & \texttt{(s)} / \texttt{(s, n)} & Reverses a string, or repeats it \texttt{n} times. \\
\textbf{LIKE / ILIKE} & pattern with \texttt{\%} and \texttt{\_} & Wildcard match, case sensitive and case insensitive respectively. \\
\textbf{\textasciitilde{} \textasciitilde* !\textasciitilde{}} & POSIX regex & Regex match, case-insensitive match, and negated match. \\
\textbf{TO\_CHAR} & \texttt{(value, format)} & Formats a number or timestamp as text, e.g.\ \texttt{'YYYY-MM-DD'}. \\
\textbf{CAST / \texttt{::}} & \texttt{(expr AS type)} / \texttt{expr::type} & Type conversion. \texttt{::numeric} also fixes integer division. \\
\bottomrule
\end{tabularx}
\caption*{String functions}
\end{table}


\newpage

% ============================================================
% DSA -- master wrapper for sections/30-code/40-dsa/.
% Pure aggregator, no heading of its own (each child keeps its
% existing \subsection). \input paths are relative to the MAIN
% document being compiled, so paths below are given in full from
% the project root.
% ============================================================

\input{sections/30-code/40-dsa/10-python-patterns}
\newpage

\input{sections/30-code/40-dsa/20-algorithmic-paradigms}
\newpage

\input{sections/30-code/40-dsa/30-recursion-memoization}
\newpage

\input{sections/30-code/40-dsa/40-heaps-priority-queues}
\newpage

\input{sections/30-code/40-dsa/50-sliding-window}

\newpage

% ============================================================
% Pandas -- master wrapper for sections/30-code/90-pandas/.
% Mirrors the sections/30-code.tex <-> sections/30-code/ pattern one
% level down. \input paths are relative to the MAIN document being
% compiled, so paths below are given in full from the project root.
% ============================================================
\subsection{Pandas}
% ============================================================

% ============================================================
\subsubsection{The Data Model}
% ============================================================

A DataFrame looks like a table, but two properties make it behave unlike a SQL table or a raw array. Every DataFrame carries a labeled axis called the index, and every column holds values of a single resolved dtype. The index drives automatic alignment across operations, and the dtype decides what an operation means and how fast it runs.

\vspace{1ex}

\begin{definition}[Series]
A Series is a one-dimensional labeled array of a single dtype. It pairs an ordered sequence of values with an \texttt{index} of labels, so every element is addressable by its label as well as by its position.
\end{definition}

\vspace{2ex}

\begin{definition}[DataFrame]
A DataFrame is a two-dimensional structure whose columns are Series sharing one common row index. Each column carries its own dtype, so a DataFrame has no single dtype of its own. The column labels form a second axis, itself an \texttt{Index}, which makes the structure symmetric: rows and columns are both labeled.
\end{definition}

A DataFrame behaves much like a dictionary of aligned Series keyed by column name. Selecting a single column returns the underlying Series, and that Series keeps the frame's row index.

\begin{lstlisting}[style=python, caption={A frame is columns of typed Series over a shared index}]
df = pd.DataFrame({
    'city': ['Austin', 'Dallas', 'Austin'],
    'temp': [98, 101, 96],
})

df['temp']        # a Series, dtype int64, indexed 0,1,2
df.index          # RangeIndex(start=0, stop=3, step=1)
df.columns        # Index(['city', 'temp'], dtype='object')
df.dtypes         # city -> object, temp -> int64
\end{lstlisting}

\vspace{-2ex}

\paragraph{The index.} A fresh DataFrame receives a default \texttt{RangeIndex} running from 0, so the labels coincide with positions and the distinction stays invisible. The index becomes load-bearing the moment it holds anything else, such as a datetime, a string key, or the grouping key that a \texttt{groupby} leaves behind. An index may be of any dtype and is not required to be unique. Two consequences shape later sections. Label-based access through \texttt{loc} and positional access through \texttt{iloc} diverge as soon as the index departs from a clean \texttt{0..n} range. And the index survives operations, so a filter or a sort carries the original labels forward rather than renumbering, which is the reason \texttt{reset\_index} exists to install a fresh range when the old labels have served their purpose.

\vspace{1ex}

\begin{keybox}[Alignment on labels]
Binary operations align their operands on the index before computing. pandas matches labels, evaluates on the ones the operands share, and fills every label present in only one operand with a missing value. Position plays no part in this matching. The same rule governs assignment and boolean-mask selection, so a mask built on one frame lands on the correct rows of another as long as the two share an index, whatever their row order.
\end{keybox}

\newpage

\begin{lstlisting}[style=python, caption={Addition aligns by label and fills the gaps}]
a = pd.Series([1, 2, 3], index=['x', 'y', 'z'])
b = pd.Series([10, 20],  index=['y', 'z'])

a + b
# x     NaN     <- 'x' has no match in b
# y    12.0
# z    23.0
\end{lstlisting}

\vspace{-1ex}

A boolean Series used as a mask, \texttt{df.loc[mask]}, is matched to the frame by label, so the mask and the frame need not share a row order for the selection to land on the right rows. The same rule cuts the other way: combining two columns taken from different frames introduces \texttt{NaN} at every label where their indices disagree, even when the columns appear to line up by position.

\paragraph{Dtypes.} Each column resolves to a single dtype, and that dtype determines both the operations available and the speed at which they run. Vectorized arithmetic and the \texttt{.str} and \texttt{.dt} accessors all require a genuine numeric, datetime, or string dtype underneath. The core dtypes are collected below.

\begin{table}[H]
\centering
\footnotesize
\renewcommand{\arraystretch}{1.4}
\rowcolors{2}{gray!10}{white}
\begin{tabularx}{\textwidth}{>{\raggedright\arraybackslash}p{0.20\textwidth} >{\raggedright\arraybackslash}p{0.28\textwidth} >{\raggedright\arraybackslash}p{0.12\textwidth} >{\raggedright\arraybackslash}X}
\toprule
\textit{dtype} & \textit{Stores} & \textit{Missing} & \textit{Note} \\
\midrule
\textbf{\texttt{int64}, \texttt{float64}} & Integers and floating point & \texttt{NaN} (float only) & The numpy-backed numeric types. \\
\textbf{\texttt{bool}} & Booleans & none & The result of every comparison. \\
\textbf{\texttt{object}} & Arbitrary Python objects, usually strings & \texttt{None}, \texttt{NaN} & Python-speed, the untyped fallback. \\
\textbf{\texttt{datetime64[ns]}} & Timestamps & \texttt{NaT} & Unlocks the \texttt{.dt} accessor. \\
\textbf{\texttt{timedelta64[ns]}} & Durations & \texttt{NaT} & Produced by subtracting timestamps. \\
\textbf{\texttt{category}} & A fixed set of repeated labels & \texttt{NaN} & Compact, and can carry an order. \\
\textbf{\texttt{Int64}, \texttt{string}, \texttt{boolean}} & As their lowercase forms & \texttt{pd.NA} & Nullable, capitalized, missing-aware. \\
\bottomrule
\end{tabularx}
\caption*{Core dtypes. The lowercase numpy-backed types are the default. The capitalized nullable types carry \texttt{pd.NA} and preserve integer and boolean columns across missing values.}
\end{table}

\vspace{-2ex}

The \texttt{object} dtype is the fallback for anything without a native type, which in practice is mostly Python strings. It stores pointers to Python objects, so work on it runs at Python speed rather than vectorized speed. Converting an \texttt{object} column of strings to a real \texttt{string} or \texttt{category} dtype, or a column of date-like text to \texttt{datetime64[ns]}, is usually the first step.

\vspace{2ex}

\begin{keybox}[Missing values force a dtype]
numpy \texttt{NaN} is a floating-point value, and numpy integer arrays have no representation for a missing entry. A single missing value in an integer column therefore promotes the whole column to \texttt{float64}. This appears constantly after a left join or a \texttt{reindex}, where unmatched rows introduce \texttt{NaN} and an \texttt{int64} count column comes back as \texttt{float64}. Datetime and timedelta columns use \texttt{NaT} in the same role. The nullable dtypes, written with a capital initial (\texttt{Int64}, \texttt{boolean}, \texttt{string}), instead carry \texttt{pd.NA} as a type-agnostic missing marker and hold an integer column at integer dtype through the gap.
\end{keybox}

\vspace{2ex}

\begin{lstlisting}[style=python, caption={A NaN forces float64, a nullable dtype keeps the integer}]
pd.Series([1, 2, 3]).dtype             # dtype('int64')
pd.Series([1, 2, np.nan]).dtype        # dtype('float64'), the gap forces the promotion

pd.Series([1, 2, pd.NA], dtype='Int64')
# 0       1
# 1       2
# 2    <NA>                            # integer dtype preserved, missing marked by pd.NA
\end{lstlisting}



% ============================================================
\subsubsection{Selection and Filtering}
% ============================================================

Three access needs recur: selecting by label, selecting by position, and selecting by predicate. pandas keeps the first two on separate accessors, \texttt{loc} for labels and \texttt{iloc} for positions, and expresses the third through a boolean mask. Restricting rows is the pandas \texttt{WHERE}, and every later pattern builds on it.

\paragraph{Bracket indexing.} The bare \texttt{[]} operator on a DataFrame is overloaded by the type of its argument. A string selects a column and returns a Series, a list of strings selects several columns and returns a DataFrame, and a boolean Series filters rows. The overloading is convenient but ambiguous, and a key that could name either a row or a column is resolved as a column. For anything past a plain column pick, \texttt{loc} states the intent without ambiguity.

\begin{lstlisting}[style=python, caption={The overloaded bracket}]
df['temp']            # a Series (one column)
df[['city', 'temp']]  # a DataFrame (a column list)
df[df['temp'] > 97]   # row filter by boolean mask
\end{lstlisting}

\vspace{-1ex}

\begin{definition}[loc and iloc]
\texttt{loc} selects by label and \texttt{iloc} by integer position. Both take a row selector and an optional column selector, \texttt{df.loc[rows, cols]}, and each selector accepts a single key, a list, a slice, or a boolean mask. A \texttt{loc} label slice includes both endpoints, whereas an \texttt{iloc} position slice excludes the stop in the usual Python way.
\end{definition}

The inclusive label slice is a frequent trap. \texttt{df.loc['a':'c']} returns rows a through c inclusive, while \texttt{df.iloc[0:3]} returns three rows and stops before position 3.

\begin{lstlisting}[style=python, caption={Label access, position access, filtered selection}]
df.loc[3, 'temp']                          # row label 3, column 'temp'
df.loc[df['temp'] > 97, ['city', 'temp']]  # mask picks rows, list picks columns
df.iloc[0]                                 # first row by position
df.iloc[-1]                                # last row by position
\end{lstlisting}

\vspace{-2ex}

\paragraph{Boolean masks.} A comparison on a Series returns a boolean Series of the same length, and passing it to \texttt{loc} keeps the \texttt{True} rows. Compound conditions combine with the bitwise operators \texttt{\&}, \texttt{|}, and \texttt{\textasciitilde}, and each comparison must be wrapped in parentheses because the bitwise operators bind tighter than the comparisons. The Python keywords \texttt{and}, \texttt{or}, and \texttt{not} do not work on a Series, since they demand a single truth value and raise on an array.

\newpage

\begin{lstlisting}[style=python, caption={Combine masks with bitwise operators, parenthesize each side}]
mask = (df['temp'] > 97) & (df['city'] == 'Austin')
df.loc[mask]

df.loc[~df['city'].isin(['Dallas', 'Houston'])]   # NOT IN
\end{lstlisting}

\vspace{-3ex}

\paragraph{Membership and range.} \texttt{isin(values)} is a vectorized membership mask, the pandas \texttt{IN}, and negating it with \texttt{\textasciitilde} gives the \texttt{NOT IN}. \texttt{between(lo, hi)} tests a closed interval. Both read more cleanly than the equivalent chain of \texttt{|} clauses or the paired inequalities they stand in for.

\paragraph{The query string.} \texttt{query} filters from a string expression, naming columns bare and outer variables with a leading \texttt{@}. It reads well for a long compound predicate and spares the repetition of the frame name. Being parsed at call time, it runs slower than a mask and accepts only expressions it can interpret, so a mask stays the general tool.

\begin{lstlisting}[style=python, caption={query names columns bare and locals with @}]
threshold = 97
df.query('temp > @threshold and city == "Austin"')
\end{lstlisting}

\begin{keybox}[Chained assignment is unreliable]
A selection can return a view onto the original's data or an independent copy, and which one arises is not predictable from the expression. Assigning through a chained selection, \texttt{df[df['t'] > 100]['city'] = 'hot'}, may therefore write into a discarded temporary and leave the original untouched, which is what raises the \texttt{SettingWithCopyWarning}. The reliable form selects the rows and the column together in a single \texttt{loc}, so the write lands on the frame itself.
\end{keybox}

\vspace{1ex}

\begin{lstlisting}[style=python, caption={The chained write warns and may no-op, the single loc is definite}]
df[df['temp'] > 100]['city'] = 'hot'          # unreliable, warns
df.loc[df['temp'] > 100, 'city'] = 'hot'      # correct, one selection
\end{lstlisting}

\vspace{-2ex}


% ============================================================
\subsubsection{Column Operations and Vectorization}
% ============================================================

pandas operations run array-at-a-time. A column expression applies to the whole Series in one pass over the underlying array, which runs far faster than a Python loop and removes the need for a loop in most transformations. This vectorized style is the default, and the row-wise \texttt{apply} is the fallback for the rare case that resists it.

\paragraph{Building columns.} Arithmetic and comparisons between columns produce new Series aligned on the index, assigned back by name. \texttt{assign} does the same inside a chain, returning a copy with columns added or overwritten. Its callable form receives the frame, which lets a new column read from the frame as it is being built.

\begin{lstlisting}[style=python, caption={Direct assignment and the chainable assign}]
df['ratio'] = df['rev'] / df['units']

df = df.assign(
    ratio=lambda d: d['rev'] / d['units'],
    is_big=lambda d: d['rev'] > 1000)
\end{lstlisting}

\begin{keybox}[Prefer vectorized operations over apply]
Arithmetic, comparisons, numpy ufuncs, and the \texttt{.str} and \texttt{.dt} accessors are all vectorized and act on the whole column at compiled speed. A Python-level \texttt{apply} runs the function once per element or row and is slower by one to two orders of magnitude. Nearly every transformation has a vectorized form, and finding it repays the effort at any nontrivial size.
\end{keybox}

\vspace{1ex}

\paragraph{Conditional columns.} A two-way choice is \texttt{np.where(cond, a, b)}, taking \texttt{a} where the condition holds and \texttt{b} elsewhere. A multi-way choice is \texttt{np.select(conditions, choices, default)}, which takes the first matching condition in order. \texttt{Series.where(cond, other)} keeps the value where the condition holds and substitutes \texttt{other} elsewhere, and \texttt{mask} is its complement. Each of these replaces a conditional \texttt{apply} with a single vectorized call.

\begin{lstlisting}[style=python, caption={Two-way and multi-way conditionals without apply}]
df['flag'] = np.where(df['temp'] > 97, 'hot', 'mild')

conditions = [df['score'] >= 90, df['score'] >= 80]
grades     = ['A', 'B']
df['grade'] = np.select(conditions, grades, default='C')
\end{lstlisting}

\vspace{-1ex}

\paragraph{The str and dt accessors.} A Series of string dtype exposes vectorized string methods under \texttt{.str}, and a datetime Series exposes calendar and clock fields under \texttt{.dt}. Both return aligned Series, so their results drop straight into masks and new columns. \texttt{pd.to\_datetime} parses text or epoch values into \texttt{datetime64[ns]}, the prerequisite for \texttt{.dt}.

\begin{lstlisting}[style=python, caption={Vectorized string and datetime work}]
df['user'].str.lower().str.startswith('a')      # chainable, vectorized
df['email'].str.contains('@company', na=False)  # na=False maps missing to False

df['ts']      = pd.to_datetime(df['ts'])
df['year']    = df['ts'].dt.year
df['weekday'] = df['ts'].dt.dayofweek           # Monday = 0
\end{lstlisting}

\vspace{-1ex}

The workhorses are \texttt{.str.contains}, \texttt{.str.extract} for a regex capture, \texttt{.str.split}, \texttt{.str.len}, and positional slicing through \texttt{.str[:3]}, alongside \texttt{.dt.year}, \texttt{.dt.month}, \texttt{.dt.date}, and \texttt{.dt.floor('D')}. The \texttt{na} argument on the string predicates fixes how a missing value scores in the resulting mask.

\paragraph{Element mapping.} \texttt{map} on a Series applies a dict or function element by element, which suits recoding a small fixed set of values. \texttt{replace} substitutes chosen values throughout a column, \texttt{astype} casts to a target dtype, and \texttt{round} sets the number of decimals on a numeric column.

\paragraph{Binning into categories.} \texttt{pd.cut} slices a numeric column at explicit edges into labeled bands, the vectorized form of a \texttt{CASE} ladder over ranges. \texttt{pd.qcut} instead cuts at quantiles, producing bins of roughly equal count, which is the \texttt{NTILE} operation. Both return a \texttt{category} column that feeds \texttt{value\_counts} or a \texttt{groupby} directly.

\newpage

\begin{lstlisting}[style=python, caption={Fixed-edge bins and equal-frequency bins}]
df['band'] = pd.cut(df['salary'],
                    bins=[0, 50_000, 80_000, np.inf],
                    labels=['<50k', '50-80k', '80k+'])

df['quartile'] = pd.qcut(df['salary'], 4, labels=False)  # 0..3 by quantile
\end{lstlisting}

\vspace{1ex}

\begin{keybox}[apply and the axis argument]
\texttt{Series.apply(f)} calls \texttt{f} once per element, and \texttt{DataFrame.apply(f, axis=1)} calls it once per row, handing each row over as a Series. Both run at Python speed and are the last resort, kept for logic with no vectorized form. The \texttt{axis} argument is a recurring source of confusion: \texttt{axis=0} runs down the index and operates on each column, while \texttt{axis=1} runs across the columns and operates on each row. A row-wise \texttt{apply} spanning several columns is the one case where \texttt{axis=1} earns its cost.
\end{keybox}


% ============================================================
\subsubsection{Grouping and Aggregation}
% ============================================================

Grouping follows the split-apply-combine shape: split the rows by a key, apply a reduction within each group, and combine the results into one row per group. It is the pandas \texttt{GROUP BY}, with the wrinkle that the grouping is a distinct lazy object holding the split, and no work happens until a reduction is named.

\vspace{1ex}

\begin{definition}[GroupBy]
\texttt{df.groupby('k')} returns a GroupBy object that partitions the rows by the values of \texttt{k} without computing anything. A reduction called on it, such as \texttt{.sum()} or \texttt{.agg(...)}, evaluates per group and combines to one row each. The key may be a list of columns for a composite grouping.
\end{definition}

\paragraph{Named aggregation.} The clearest reduction form names each output column and gives its source column and function as a tuple, \texttt{out=('col', 'func')}. It maps onto \texttt{func(col) AS out} directly, and it extends to several outputs in one call.

\begin{lstlisting}[style=python, caption={Named aggregation, one row per group}]
g = (df.groupby('customer_id', as_index=False)
       .agg(n_orders=('order_id', 'count'),
            total=('amount', 'sum')))
\end{lstlisting}

\vspace{-1ex}

The older forms remain available. A dict \texttt{agg(\{'amount': 'sum'\})} maps a column to a function, and a list \texttt{agg(['sum', 'mean'])} applies several functions at once and returns a frame with a MultiIndex on its columns. The named form sidesteps that nested column index.

\paragraph{The group key becomes the index.} A \texttt{groupby} puts the key on the index of the result, one level per key column. \texttt{reset\_index()} moves it back to an ordinary column and installs a clean range index, and passing \texttt{as\_index=False} to \texttt{groupby} does the same inline. The key needs to leave the index whenever a later step selects it as a column or expects a plain integer row label.

\vspace{1ex}

\begin{keybox}[count versus size versus nunique]
\texttt{count} tallies non-null values in the group, \texttt{size} tallies rows including nulls, and \texttt{nunique} tallies distinct non-null values. The gap between \texttt{count} and \texttt{size} is where left-join placeholders hide: an unmatched row carries \texttt{NaN}, so \texttt{size} scores it as one row while \texttt{count} scores it as zero. Choosing the intended one of these is a routine source of off-by-$N$ errors.
\end{keybox}

\vspace{2ex}

\begin{table}[H]
\centering
\footnotesize
\renewcommand{\arraystretch}{1.4}
\rowcolors{2}{gray!10}{white}
\begin{tabularx}{\textwidth}{>{\raggedright\arraybackslash}p{0.15\textwidth} >{\raggedright\arraybackslash}p{0.30\textwidth} >{\raggedright\arraybackslash}X}
\toprule
\textit{Reduction} & \textit{SQL analog} & \textit{What it computes} \\
\midrule
\textbf{\texttt{count}} & \texttt{COUNT(col)} & Non-null values in the group. \\
\textbf{\texttt{size}} & \texttt{COUNT(*)} & Every row in the group, nulls included. \\
\textbf{\texttt{nunique}} & \texttt{COUNT(DISTINCT col)} & Distinct values, nulls excluded. \\
\textbf{\texttt{sum}} & \texttt{SUM(col)} & Total of the values. An empty group sums to 0. \\
\textbf{\texttt{mean}} & \texttt{AVG(col)} & Arithmetic mean. \\
\textbf{\texttt{median}} & \texttt{PERCENTILE\_CONT(0.5)} & Middle value by order. \\
\textbf{\texttt{min}} & \texttt{MIN(col)} & Smallest value. \\
\textbf{\texttt{max}} & \texttt{MAX(col)} & Largest value. \\
\textbf{\texttt{std}} & \texttt{STDDEV\_SAMP(col)} & Sample standard deviation, \texttt{ddof = 1}. \\
\textbf{\texttt{var}} & \texttt{VAR\_SAMP(col)} & Sample variance, \texttt{ddof = 1}. \\
\textbf{\texttt{first}} & \texttt{FIRST\_VALUE} (window) & First value in the current row order. \\
\textbf{\texttt{last}} & \texttt{LAST\_VALUE} (window) & Last value in the current row order. \\
\textbf{\texttt{any / all}} & \texttt{BOOL\_OR / BOOL\_AND} & Whether any, or all, values are truthy. \\
\bottomrule
\end{tabularx}
\caption*{Common reductions for \texttt{agg}. All skip \texttt{NaN} except \texttt{size}. \texttt{std} and \texttt{var} default to the sample estimator, unlike SQL engines that default to population.}
\end{table}

\vspace{-2ex}

\paragraph{transform versus agg.} \texttt{agg} returns one row per group. \texttt{transform} returns a result the same length as the input, broadcasting each group's statistic back onto every member row. It is the tool for a group-relative column, where each row is measured against a summary of its own group.

\begin{lstlisting}[style=python, caption={A group-relative column via transform}]
dept_mean = df.groupby('dept')['salary'].transform('mean')
df['above_dept_avg'] = df['salary'] > dept_mean
\end{lstlisting}

\vspace{-1ex}

\paragraph{Filtering whole groups.} \texttt{groupby(...).filter(pred)} keeps or drops entire groups by a group-level predicate, the reading of \texttt{HAVING}. The predicate runs once per group in Python, so on a wide key the vectorized equivalent through \texttt{transform('size')} and a mask runs faster.

\begin{lstlisting}[style=python, caption={Keep groups by a group predicate (HAVING)}]
# groups of size greater than one
df.groupby('email').filter(lambda g: len(g) > 1)

# vectorized equivalent
df[df.groupby('email')['email'].transform('size') > 1]
\end{lstlisting}

\vspace{2ex}

\begin{keybox}[groupby drops null keys by default]
A row whose grouping key is \texttt{NaN} is excluded from every group, so a null-keyed record drops out of the result. Passing \texttt{dropna=False} retains the null-key group. Set it whenever a grouping key may be missing and those rows still belong in the output.
\end{keybox}

\vspace{-2ex}

\paragraph{Frequencies and distinct values.} \texttt{value\_counts} on a Series tallies how often each value occurs and returns the counts sorted high to low, with \texttt{normalize=True} giving proportions instead. \texttt{unique} lists the distinct values in order of appearance and \texttt{nunique} counts them. \texttt{drop\_duplicates} removes repeated rows, keyed to chosen columns through \texttt{subset} and keeping the first or last of each run through \texttt{keep}. A grouped string concatenation, the \texttt{GROUP\_CONCAT}, falls out of \texttt{agg} with a join function.

\begin{lstlisting}[style=python, caption={Frequencies, distinct rows, and joined strings}]
df['dept'].value_counts(normalize=True)      # share of rows per dept
df.drop_duplicates(subset=['email'], keep='first')

df.groupby('dept')['name'].agg(', '.join)    # GROUP_CONCAT
\end{lstlisting}


% ============================================================
\subsubsection{Joining and Combining}
% ============================================================

Two operations bring separate frames together. \texttt{merge} combines columns by matching keys, the pandas \texttt{JOIN}. \texttt{concat} stacks frames along an axis without matching. Reach for \texttt{merge} when rows are paired by a key and for \texttt{concat} when frames are laid end to end.

\vspace{1ex}

\begin{definition}[merge]
\texttt{merge} joins two frames on one or more key columns. \texttt{on} names a key shared by both, while \texttt{left\_on} and \texttt{right\_on} pair keys that differ in name. \texttt{how} selects which unmatched rows survive, taking \texttt{inner}, \texttt{left}, \texttt{right}, \texttt{outer}, or \texttt{cross}, and defaults to \texttt{inner}.
\end{definition}

The \texttt{how} argument carries the same meaning as the SQL join it names. \texttt{inner} keeps rows matched on both sides. \texttt{left} and \texttt{right} keep every row from one side and fill the other with \texttt{NaN}. \texttt{outer} keeps unmatched rows from both. \texttt{cross} forms the Cartesian product and takes no key.

\begin{lstlisting}[style=python, caption={A keyed merge and a rename-free key pairing}]
orders.merge(customers, on='customer_id', how='left')
a.merge(b, left_on='cid', right_on='id', how='inner')
\end{lstlisting}

\begin{keybox}[Join type decides which unmatched rows survive]
The \texttt{how} argument decides whether an unmatched row survives as a hole or disappears. A left join keeps every left row and fills the unmatched right columns with \texttt{NaN}, while an inner join drops those rows outright. The choice carries a dtype consequence: filling an integer column with \texttt{NaN} promotes it to \texttt{float64}, so an \texttt{int64} count returns as \texttt{float64} after a left join.
\end{keybox}

\vspace{1ex}

\paragraph{Overlapping column names.} When both frames carry a non-key column of the same name, \texttt{merge} disambiguates the two with the default suffixes \texttt{\_x} and \texttt{\_y}. The result then requires tracking which side is which, so passing an explicit \texttt{suffixes=} pair, or renaming beforehand, keeps the output readable.

\paragraph{Joining on the index.} \texttt{join} is a convenience wrapper over \texttt{merge} that matches on the index of the right frame by default, which suits attaching a lookup table already indexed by its key. \texttt{merge} stays the explicit tool for arbitrary key columns, and it reads more clearly whenever the key is an ordinary column rather than the index.

\vspace{2ex}

\begin{keybox}[Duplicate keys multiply rows]
A merge emits one row per matching pair of keys. When a key repeats on both sides, the matches fan out as their product, so a many-to-many merge can return far more rows than either input. Passing \texttt{validate='one\_to\_one'} or \texttt{'one\_to\_many'} makes \texttt{merge} assert the expected multiplicity and raise on a violation.
\end{keybox}

\paragraph{Stacking with concat.} \texttt{concat} joins a list of frames along one axis. \texttt{axis=0} stacks rows and aligns on the columns, \texttt{axis=1} stacks columns and aligns on the index. \texttt{ignore\_index=True} discards the source indices and lays down a fresh range index, which matters when the pieces carry overlapping or meaningless labels.

\begin{lstlisting}[style=python, caption={Stack rows, then stack columns}]
pd.concat([jan, feb, mar], ignore_index=True)   # rows, fresh index
pd.concat([features, labels], axis=1)           # columns, aligned on index
\end{lstlisting}

\vspace{-1ex}

\paragraph{Semi- and anti-joins through isin.} Filtering a frame by whether its key appears in another needs no materialized join. \texttt{isin} against the other frame's key column keeps the matching rows, a semi-join, and negating it keeps the non-matching rows, an anti-join. Neither widens the frame nor pulls columns from the other side.

\begin{lstlisting}[style=python, caption={Semi-join keeps matches, anti-join keeps non-matches}]
customers[customers['id'].isin(orders['customer_id'])]    # has orders
customers[~customers['id'].isin(orders['customer_id'])]   # has no orders
\end{lstlisting}

\begin{keybox}[Filter anti-joins on a key]
An exclusion should test the column that identifies a row uniquely, which is the primary key. Filtering on a display attribute such as a name folds together distinct records that happen to share it, so one person's match can drop an unrelated person of the same name from the result. The membership form also holds an advantage over the SQL \texttt{NOT IN}: \texttt{isin} returns an independent boolean per row, so a missing value among the excluded keys cannot void the whole result the way a single \texttt{NULL} collapses \texttt{NOT IN} to zero rows.
\end{keybox}

\paragraph{As-of joins.} \texttt{merge\_asof} matches each left row to the single right row whose key is nearest without exceeding it, the backward direction being the default. It answers point-in-time questions, where a fact is read as of a moment rather than on an exact match. Both frames must be sorted on the join key, since the algorithm advances through them in order. The \texttt{by} argument confines matching within a group, pairing a row only with keys from the same group.

\newpage

\begin{lstlisting}[style=python, caption={Match each sale to the price in effect at its date}]
prices = prices.sort_values('start_date')
sales  = sales.sort_values('purchase_date')

pd.merge_asof(sales, prices, by='product_id',
              left_on='purchase_date', right_on='start_date')
\end{lstlisting}

\vspace{-1ex}

The join consults the start key alone and never a matching end key, so it is not a between-dates join and will attach a stale value to a left row falling past a window's close. It fits when each key marks the start of an interval running until the next one.


% ============================================================
\subsubsection{Reshaping}
% ============================================================

Reshaping moves data between the row axis and the column axis. In long form each row is one observation, keyed by its identifying columns. In wide form the values of one key spread across the columns, one column per distinct value. \texttt{pivot} carries a frame from long to wide, and \texttt{melt} carries it back.

\vspace{1ex}

\begin{definition}[pivot and pivot\_table]
\texttt{pivot(index, columns, values)} lifts the distinct entries of the \texttt{columns} field onto the column axis, indexes the result by \texttt{index}, and fills each cell from \texttt{values}. It requires exactly one value per \texttt{(index, columns)} cell and raises on a duplicate. \texttt{pivot\_table} adds an \texttt{aggfunc} that reduces colliding values, so it tolerates duplicates by aggregating them.
\end{definition}

\paragraph{Pivoting to pair split rows.} A single logical record is sometimes spread across several rows, distinguished by a type column with one measured value each. Pivoting that type column lifts its values into columns, bringing the pieces of one record onto a single row where they combine through column arithmetic. This replaces a self-join for the common case of a small fixed set of types, such as a start and an end.

\begin{lstlisting}[style=python, caption={Pivot a start/end type column into aligned columns}]
wide = events.pivot(index=['machine_id', 'process_id'],
                    columns='activity_type',
                    values='timestamp')

wide['duration'] = wide['end'] - wide['start']
\end{lstlisting}

\begin{keybox}[pivot versus pivot\_table]
\texttt{pivot} reshapes without aggregating and raises when a cell would receive more than one value. \texttt{pivot\_table} aggregates collisions through its \texttt{aggfunc}, which defaults to the mean. Use \texttt{pivot} when uniqueness is guaranteed by the data, so an unexpected duplicate surfaces as an error rather than being averaged away. Reach for \texttt{pivot\_table} when collisions are expected and a summary of them is the intent.
\end{keybox}

\paragraph{Melting wide to long.} \texttt{melt} is the inverse operation. The \texttt{id\_vars} columns stay fixed as identifiers, and the \texttt{value\_vars} columns collapse into two columns, one holding the former column name and one holding its value. It normalizes a wide frame into the long form that grouping and plotting expect.

\begin{lstlisting}[style=python, caption={Collapse quarter columns into rows}]
long = wide.melt(id_vars='region',
                 value_vars=['q1', 'q2', 'q3', 'q4'],
                 var_name='quarter', value_name='revenue')
\end{lstlisting}

\vspace{-2ex}

\paragraph{stack and unstack.} \texttt{stack} moves a column level down into a row-index level, and \texttt{unstack} lifts a row-index level up into columns. They are the index-level counterparts of \texttt{melt} and \texttt{pivot}, and \texttt{pivot} itself is \texttt{set\_index} followed by \texttt{unstack}. Reshaping often leaves a MultiIndex on the rows or columns, which \texttt{reset\_index}, \texttt{droplevel}, or \texttt{rename\_axis} flatten back to plain labels.


% ============================================================
\subsubsection{Ordering and Row-Relative Logic}
% ============================================================

Some results depend on the order of the rows or on values in other rows. Ordering operations sort or rank the frame, and row-relative operations read a neighbor or a running window. These are the pandas counterpart of SQL window functions, and as with a SQL window the order must be stated, since a frame carries no inherent one.

\paragraph{Sorting.} \texttt{sort\_values(by, ascending, na\_position)} orders rows by one or more columns, with \texttt{by} taking a list for a multi-key sort and \texttt{ascending} a matching list of directions. \texttt{na\_position} places \texttt{NaN} at the start or the end. \texttt{sort\_index} orders by the index labels instead.

\begin{lstlisting}[style=python, caption={Multi-key sort with a direction per key}]
df.sort_values(['dept', 'salary'], ascending=[True, False])
\end{lstlisting}

\vspace{-1ex}

\paragraph{Extremes and top-N.} \texttt{nlargest(n, col)} and \texttt{nsmallest(n, col)} return the top or bottom rows by a column and cost less than a full sort for a handful of rows. \texttt{idxmax} and \texttt{idxmin} return the index label of the extreme value, an argmax, which \texttt{loc} then reads as a row. Both break ties by first occurrence.

\begin{lstlisting}[style=python, caption={Top rows and the row of the maximum}]
df.nlargest(3, 'salary')
df.loc[df['salary'].idxmax()]      # the single row holding the max
\end{lstlisting}

\vspace{-1ex}

\paragraph{The Nth distinct value.} The second-highest and Nth-highest questions read the value sitting at a chosen rank. Deduplicating the values first keeps ties from consuming a rank, and the result resolves to a missing value when the data holds fewer than \texttt{N} distinct entries, which is the required answer to an absent Nth place.

\begin{lstlisting}[style=python, caption={Nth-highest distinct value, NaN when it does not exist}]
def nth_highest(s, n):
    u = s.drop_duplicates().nlargest(n)
    return u.iloc[-1] if len(u) >= n else np.nan
\end{lstlisting}

\vspace{-1ex}

\paragraph{Ranking.} \texttt{rank(method, ascending)} numbers rows by value, with the tie policy set by \texttt{method}: \texttt{dense}, \texttt{min}, and \texttt{first} correspond to the SQL \texttt{DENSE\_RANK}, \texttt{RANK}, and \texttt{ROW\_NUMBER}. Applied after a \texttt{groupby} it ranks within each group, the top-$N$-per-group tool. Sorting and then \texttt{groupby(...).head(n)} reaches the same result by keeping the first \texttt{n} rows of each group. Passing \texttt{pct=True} scales the ranks to the interval $(0, 1]$, giving a percent rank.

\begin{lstlisting}[style=python, caption={Rank within a group, and top-N per group by slicing}]
df['rk'] = df.groupby('dept')['salary'].rank(method='dense', ascending=False)

(df.sort_values('salary', ascending=False)
   .groupby('dept').head(1))        # highest-paid row per dept
\end{lstlisting}

\vspace{-1ex}

\paragraph{Position within a group.} \texttt{groupby(...).cumcount()} numbers the rows of each group by their current order from 0 and resets at every group edge, the within-group \texttt{ROW\_NUMBER} carrying no value to rank on. \texttt{ngroup()} labels the groups themselves \texttt{0, 1, 2} in order of appearance. Neither reads a column, so both rest on a prior sort the way \texttt{rank} does.

\begin{lstlisting}[style=python, caption={A 0-based sequence number within each group}]
df['seq'] = df.sort_values('ts').groupby('user').cumcount()
\end{lstlisting}

\vspace{-1ex}

\begin{keybox}[Row-relative operations need an explicit order]
\texttt{shift}, \texttt{diff}, \texttt{rank}, \texttt{rolling}, and the cumulative methods all read the current row order and nothing else. Without a preceding \texttt{sort\_values} the previous row is whatever order the rows arrived in. This is the \texttt{ORDER BY} inside a SQL window, and a per-group version needs the operation applied through \texttt{groupby} so it resets at each group boundary rather than bleeding across.
\end{keybox}

\vspace{-1ex}

\paragraph{Neighbor access with shift and diff.} \texttt{shift(n)} moves a column down \texttt{n} positions and fills the vacated top with \texttt{NaN}, so \texttt{shift(1)} sets the previous row's value beside the current one, the \texttt{LAG}. \texttt{shift(-1)} looks ahead, the \texttt{LEAD}. \texttt{diff()} equals \texttt{current - shift(1)}, so \texttt{diff() > 0} tests whether a value rose. The first row of an ordering has no predecessor, so its shifted value is \texttt{NaN} and every comparison against it is \texttt{False}, which drops it with no explicit guard. Applying the shift through \texttt{groupby} resets it at each group edge. The one-period fractional change, \texttt{close / close.shift(1) - 1}, is packaged as \texttt{pct\_change}, the simple return.

\begin{lstlisting}[style=python, caption={Per-group return against the previous row}]
df = df.sort_values(['ticker', 'date'])

prev = df.groupby('ticker')['close'].shift(1)
df['ret'] = df['close'] / prev - 1               # explicit, via shift
df['ret'] = df.groupby('ticker')['close'].pct_change()   # same return
\end{lstlisting}

\vspace{-2.5ex}

\paragraph{Running and rolling windows.} A cumulative method such as \texttt{cumsum} or \texttt{cummax} runs from the start up to the current row, giving a running total. \texttt{rolling(k)} fixes a window of \texttt{k} consecutive rows for a moving statistic, and \texttt{expanding} grows the window from the start. \texttt{min\_periods} sets how many observations a window needs before it returns a value in place of \texttt{NaN}. Wrapping these in \texttt{groupby(...).transform(...)} keeps the result aligned to the original rows while resetting each group.

\begin{lstlisting}[style=python, caption={A per-group moving average aligned to the rows}]
df['ma20'] = (df.groupby('ticker')['close']
                .transform(lambda s: s.rolling(20).mean()))
\end{lstlisting}

\vspace{-1ex}

\texttt{rolling} accepts any reduction, so \texttt{.rolling(20).std()} is a moving volatility, and a moving correlation between two aligned series would be \texttt{.rolling(60).corr(other)}. \texttt{ewm} weights recent observations more heavily under an exponential decay set by \texttt{span}, \texttt{halflife}, or \texttt{alpha}, giving the exponentially weighted mean and variance used in volatility estimates. Compounding runs through \texttt{cumprod}, so \texttt{(1 + ret).cumprod() - 1} accumulates a return series.

\begin{lstlisting}[style=python, caption={Moving volatility, exponential mean, and compounded return}]
g = df.groupby('ticker')
df['vol20']   = g['ret'].transform(lambda s: s.rolling(20).std())
df['ewma20']  = g['close'].transform(lambda s: s.ewm(span=20).mean())
df['cum_ret'] = g['ret'].transform(lambda s: (1 + s).cumprod() - 1)
\end{lstlisting}

\vspace{-1ex}

\paragraph{Time-based and custom windows.} An integer \texttt{rolling(k)} counts \texttt{k} rows and ignores their timestamps. An offset string, \texttt{rolling('7D')}, counts a span of calendar time instead and needs a \texttt{DatetimeIndex} or an \texttt{on=} column naming the time field. The two coincide only when the rows sit one period apart with no gap between them. A window takes any reduction, so \texttt{sum}, \texttt{min}, \texttt{max}, \texttt{median}, and \texttt{quantile} all apply, alongside \texttt{center=True} for a window centered on each row. \texttt{rolling(k).apply(f, raw=True)} runs an arbitrary \texttt{f} over each window as a numpy array, the route for a statistic with no built-in form.

\begin{lstlisting}[style=python, caption={A calendar window, a custom window, and a centered median}]
df = df.set_index('date').sort_index()

df['roll7'] = df['amount'].rolling('7D').sum()     # seven calendar days back
df['rng10'] = df['amount'].rolling(10).apply(lambda w: w.max() - w.min(), raw=True)
df['med5']  = df['amount'].rolling(5, center=True).median()
\end{lstlisting}

\vspace{2ex}

\begin{keybox}[A row window is not a day window]
\texttt{rolling(7)} spans seven rows, which equals seven days only when every day is present. On a series with missing dates it reaches further back in time than seven days. Two fixes hold the window to calendar length: reindex to a complete date range so each day is a row, or use the offset form \texttt{rolling('7D')}, which measures the span directly and covers the rows whose timestamps fall within seven days back of the current one.
\end{keybox}

\vspace{-1ex}

\paragraph{Time indexing and resampling.} A \texttt{DatetimeIndex} unlocks time-aware operations the ordinary index lacks. Partial-string slicing selects by calendar span, so \texttt{df.loc['2023']} takes the whole year and \texttt{df.loc['2023-01':'2023-06']} a range of months. \texttt{resample} is a groupby keyed by a time frequency: it bins the rows into calendar buckets and reduces each, downsampling a series to a coarser frequency or upsampling to a finer one that then needs filling. Frequency aliases name the bucket, among them \texttt{'B'} business day, \texttt{'W'} week, \texttt{'ME'} month-end, and \texttt{'QE'} quarter-end, the last two written \texttt{'M'} and \texttt{'Q'} before pandas 2.2.

\begin{lstlisting}[style=python, caption={Downsample, build OHLC bars, and align to business days}]
px = px.set_index('date').sort_index()

px['close'].resample('ME').last()      # month-end close (downsample)
px['close'].resample('W').ohlc()       # weekly open-high-low-close bars
px['close'].resample('B').ffill()      # to business days, carry last forward

px.groupby('ticker')['close'].resample('ME').last()   # per entity
\end{lstlisting}


\newpage

\begin{center}
{\Large {\color{red} Style Reference Page}}
\end{center}
\vspace{5em}

\begin{definition}[Expectation]
For a discrete random variable $X$ with support $\mathcal{X}$,
$\E[X] = \sum_{x \in \mathcal{X}} x \, \Prob(X = x)$, provided the sum converges absolutely.
\end{definition}

\begin{keybox}[Linearity of expectation]
For any random variables $X_1,\dots,X_n$ and scalars $a_i$,
\[
  \E\!\left[\sum_{i=1}^n a_i X_i\right] = \sum_{i=1}^n a_i \, \E[X_i].
\]
This holds with no independence assumption. It is the single highest-yield identity for interview brainteasers, especially when combined with indicator variables $\ind\{\cdot\}$.
\end{keybox}

\begin{theorem}[Variance of a sum]
If $X_1,\dots,X_n$ are pairwise uncorrelated, then
$\Var\!\left(\sum_i X_i\right) = \sum_i \Var(X_i)$.
\end{theorem}

\begin{proof}
Expand $\Var(\sum_i X_i) = \sum_i \Var(X_i) + \sum_{i \neq j}\Cov(X_i,X_j)$.
Pairwise uncorrelatedness kills every cross term, leaving the claim.
\end{proof}

\begin{example}
Draw a uniformly random permutation of $\{1,\dots,n\}$. Let $X$ count fixed points.
Write $X = \sum_i \ind\{\pi(i)=i\}$. Each indicator has mean $1/n$, so by linearity $\E[X]=1$, independent of $n$.
\todo{Add the variance derivation and the derangement limit $\Prob(X=0)\to e^{-1}$.}
\end{example}

\begin{remark}
Keep a running list of distributions with their mean, variance, and MGF. Interviewers probe recall speed on Bernoulli, Binomial, Geometric, Poisson, Exponential, and Normal.
\end{remark}


%

\end{document}
